% =====================================================================
% Applied Finance — Working Paper n. 1
% Cinco abordagens de seleção de carteira no mercado brasileiro:
% replicação de Markowitz, DGU, Kahneman, Black-Litterman e Paridade
% de Risco sobre o universo IBOV via lakehouse Databricks.
%
% Autor: Leonardo Chalhoub
% Versão: 1.0 — Maio de 2026
% =====================================================================
\documentclass[12pt,a4paper,oneside]{article}

\usepackage[T1]{fontenc}
\usepackage[utf8]{inputenc}
\usepackage[brazilian]{babel}
\usepackage{newtxtext}
\usepackage{newtxmath}
\usepackage[section]{placeins}
\usepackage[a4paper,top=3cm,bottom=2cm,left=3cm,right=2cm]{geometry}
\usepackage{setspace}
\usepackage{indentfirst}
\usepackage{titlesec}
\usepackage{caption}
\usepackage{booktabs}
\usepackage{array}
\usepackage{multirow}
\usepackage{graphicx}
\usepackage[colorlinks=true,
            urlcolor=blue!60!black,
            linkcolor=black,
            citecolor=black,
            breaklinks=true]{hyperref}
\usepackage{url}
\usepackage{xurl}
\Urlmuskip=0mu plus 1mu\relax
\usepackage{float}
\usepackage{amsmath}
% Note: amssymb intentionally NOT loaded — newtxmath already provides
% the full AMS symbol set, and loading amssymb after newtxmath triggers
% "Command \Bbbk already defined" in newer TeX Live (2026+).

\onehalfspacing
\setlength{\parindent}{1.25cm}
\setlength{\parskip}{0pt}

\titleformat{\section}{\normalfont\bfseries\large\uppercase}{\thesection}{0.6em}{}
\titleformat{\subsection}{\normalfont\bfseries\normalsize}{\thesubsection}{0.6em}{}
\titlespacing*{\section}{0pt}{18pt}{8pt}
\titlespacing*{\subsection}{0pt}{12pt}{6pt}

\captionsetup[table]{labelfont=bf,name=Tabela,labelsep=endash,
                     position=above,justification=centering,font=small}
\captionsetup[figure]{labelfont=bf,name=Figura,labelsep=endash,
                      position=below,justification=centering,font=small}
\captionsetup{skip=6pt}

\newcommand{\rs}{R\$\,}

% Auto-generated in CI by .github/workflows/deploy-pages.yml — do NOT edit by hand.
% Locally this default file is used and the cover page shows "(compilação local)".
\newcommand{\COMPILESTAMP}{compila\c{c}\~ao local}
\newcommand{\COMPILESHA}{dev}


\begin{document}

% ---------- CAPA ---------------------------------------------------------
\begin{titlepage}
  \begin{center}
    \vspace*{1.5cm}
    {\small\textsc{Applied Finance \quad\textbar\quad Working Paper n. 1}}\\
    {\small Versão 1.0 \textemdash{} Maio de 2026}\\[2pt]
    {\footnotesize Compilada em \COMPILESTAMP\ \textbullet\ \texttt{\COMPILESHA}}

    \vfill

    {\LARGE\bfseries
     CINCO ABORDAGENS DE SELEÇÃO DE CARTEIRA NO MERCADO BRASILEIRO:

     REPLICAÇÃO EMPÍRICA DE MARKOWITZ, DEMIGUEL-GARLAPPI-UPPAL,

     KAHNEMAN, BLACK-LITTERMAN E PARIDADE DE RISCO

     SOBRE O UNIVERSO IBOV (2000\textendash 2026)}

    \vspace{1cm}
    {\large\itshape
     Five portfolio-selection approaches on the Brazilian equity market:
     empirical replication of Markowitz, DeMiguel-Garlappi-Uppal, Kahneman,
     Black-Litterman and Risk Parity on the IBOV universe (2000\textendash 2026)}

    \vfill

    {\large\textbf{Leonardo Chalhoub}}\\[4pt]
    {\small Pesquisador independente \textemdash{} engenharia de dados,
     finanças quantitativas e arquitetura \textit{lakehouse}.}\\[2pt]
    {\small\href{mailto:leonardochalhoub@gmail.com}{leonardochalhoub@gmail.com}}\\
    {\small\href{https://github.com/leonardochalhoub/applied\_finance}{github.com/leonardochalhoub/applied\_finance}}

    \vfill

    {\small Brasil \\ Maio de 2026}
  \end{center}
\end{titlepage}

% ---------- RESUMO -------------------------------------------------------
\section*{Resumo}
\addcontentsline{toc}{section}{Resumo}
\begin{singlespace}\small\noindent
A literatura quantitativa de seleção de carteira oscila entre dois pólos:
o ideal teórico de Markowitz (1952), que prescreve a otimização de
média-variância como solução única e completa do problema do investidor
sob preferências quadráticas; e a crítica empírica acumulada desde
Michaud (1989) e formalizada por DeMiguel, Garlappi e Uppal (2009), que
documenta que a média amostral $\hat{\mu}$ é estatisticamente frágil
demais para que o otimizador supere consistentemente regras ingênuas
fora-da-amostra. Este \textit{working paper} replica empiricamente cinco
abordagens contemporâneas \textemdash{} Markowitz puro, ingênuo $1/N$
de DeMiguel-Garlappi-Uppal, \textit{Illusion of Skill} de Kahneman (2011,
prêmio Nobel de Ciências Econômicas 2002), Black-Litterman (1992) e
Paridade de Risco / Equal Risk Contribution de Maillard, Roncalli e
Teïletche (2010) \textemdash{} sobre o universo de ações listadas na
B3, em janela 2000\textendash 2026.
\textbf{Materiais e métodos.} A fonte de dados é uma plataforma
\textit{lakehouse} Medallion (\textit{bronze} $\rightarrow$ \textit{silver}
$\rightarrow$ \textit{gold}) sobre Databricks Free Edition, com 982
\textit{tickers} no universo curado, 468 efetivamente ingeridos pelo
\textit{provider} (Yahoo Finance), 38 constituintes correntes do IBOV
disponíveis em série de preços fechados ajustados. Estimadores: matriz
de covariância via Ledoit-Wolf (2004); média esperada via correção de
Jensen, anualização $\times 252$, encolhimento Bayes-Stein de Jorion
(1986) e ancoragem macro a $r_f + \mathrm{ERP}$ com teto por ativo em
$r_f + 3 \cdot \mathrm{ERP} \approx 31\%$; solver \textit{long-only}
via heurística \textit{active-set} guloso.
\textbf{Achados.} (i)~A fronteira Markowitz \textit{long-only} sobre as
30 maiores constituintes do IBOV ($T = 1260$ dias, $r_f = 12{,}90\%$ a.a.)
entrega tangência com $E[r] = 28{,}52\%$, $\sigma = 21{,}56\%$ e
Sharpe $= 0{,}72$, concentrando $26\% + 25\% + 18\% = 69\%$ em
SBSP3, PETR4 e PETR3. (ii)~Em \textit{walk-forward} 3-anos-treino /
1-trimestre-teste sobre 8 períodos entre abril/2024 e maio/2026, a
carteira Markowitz supera a $1/N$ em retorno acumulado
($+59{,}65\%$ vs $+39{,}17\%$) e em Sharpe anualizado
($0{,}73$ vs $0{,}27$), com \textit{turnover} de $1{,}17\times$ ao
ano \textemdash{} um resultado que \textit{contradiz} a tese DGU no
intervalo considerado mas reflete uma janela de mercado favorável aos
ativos sobreponderados pela otimização. (iii)~No experimento de
Kahneman (5.000 carteiras Dirichlet long-only, partição treino/teste
$2y$/$2y$), o Sharpe ex-ante prometido pelo otimizador é $1{,}123$
contra $0{,}947$ ex-post; o gap, $0{,}176$ em unidades de Sharpe,
constitui a \textit{ilusão} no sentido formal de Kan e Smith (2008);
ambos valores ficam acima da totalidade do suporte aleatório, indicando
\textit{skill real} estatisticamente distinguível de sorte na janela
considerada. (iv)~O modelo Black-Litterman, com $\delta = 2{,}5$ e
$\tau = 0{,}05$ e sem \textit{views}, recupera implícitos
$\Pi$ entre $5{,}64\%$ (ABEV3) e $10{,}49\%$ (B3SA3), distantes do
$\hat{\mu}$ amostral encolhido pelo pipeline ($12{,}64\%$ a $30{,}90\%$).
A norma $L^1/2$ do desvio em relação à carteira de mercado é $70{,}5\%$
para Markowitz e $72{,}4\%$ para Black-Litterman. (v)~A carteira de
paridade de risco (ERC) atinge $\mathrm{HHI}_{\text{risco}} = 0{,}067
\approx 1/15$ por construção, contra $0{,}069$ no $1/N$ e $0{,}318$
em Markowitz; a vol total é $17{,}84\%$, $18{,}35\%$ e $22{,}99\%$
respectivamente. \textbf{Discussão e contribuição.} O \textit{paper}
reúne os cinco modelos sob um \textit{único} pipeline de estimação e
contraste empírico, com pesos comparáveis em base ativo-a-ativo; o
mesmo $\Sigma$ Ledoit-Wolf é usado em todas as etapas, garantindo
comparabilidade interna. A reprodução é integral: os algoritmos são
\textit{open-source} (TypeScript, com \textit{tests} pytest e vitest),
os dados são públicos (Yahoo Finance + Banco Central do Brasil), e a
plataforma é gratuita (Databricks Free Edition).
\vspace{0.4cm}\noindent\textbf{Palavras-chave:}
Markowitz; DeMiguel-Garlappi-Uppal; Kahneman; Black-Litterman; paridade
de risco; \textit{Ledoit-Wolf}; Jorion; carteira de tangência;
\textit{lakehouse} Medallion; Databricks; B3; IBOV; reprodutibilidade.
\end{singlespace}

% ---------- ABSTRACT -----------------------------------------------------
\section*{Abstract}
\addcontentsline{toc}{section}{Abstract}
\begin{singlespace}\small\noindent
The quantitative literature on portfolio selection oscillates between
two poles: Markowitz's (1952) theoretical ideal, prescribing mean-variance
optimization as the unique and complete solution to the investor's
problem under quadratic preferences; and the empirical critique
accumulated since Michaud (1989), formalized by DeMiguel, Garlappi and
Uppal (2009), documenting that the sample mean $\hat{\mu}$ is too
statistically fragile for the optimizer to consistently outperform
naïve rules out-of-sample. This working paper empirically replicates
five contemporary approaches \textemdash{} pure Markowitz, naïve
$1/N$, Kahneman's \textit{Illusion of Skill} (2011, Nobel 2002),
Black-Litterman (1992) and Equal Risk Contribution (Maillard, Roncalli
\& Teïletche, 2010) \textemdash{} on the B3 equity universe, window
2000\textendash 2026.
\textbf{Methods.} Data come from a Medallion lakehouse
(\textit{bronze} $\rightarrow$ \textit{silver} $\rightarrow$
\textit{gold}) on Databricks Free Edition, with 982 curated tickers,
468 effectively ingested via Yahoo Finance, 38 current IBOV
constituents with adjusted-close series. Estimators: Ledoit-Wolf
covariance, Jensen-corrected $\times 252$ annualized mean, Jorion
(1986) Bayes-Stein shrinkage and macro-anchor to $r_f + \mathrm{ERP}$
with per-asset ceiling at $r_f + 3 \cdot \mathrm{ERP} \approx 31\%$;
long-only solver via greedy active-set heuristic.
\textbf{Findings.} The long-only Markowitz tangency on top-30 IBOV
($T = 1260$, $r_f = 12.90\%$ p.a.) delivers $E[r] = 28.52\%$,
$\sigma = 21.56\%$, Sharpe $= 0.72$. Walk-forward over 8 periods
(Apr/2024\textendash May/2026) shows Markowitz beating $1/N$ in this
specific bull window (Sharpe $0.73$ vs $0.27$); the Kahneman experiment
records ex-ante Sharpe $1.12$ vs ex-post $0.95$, illusion $+0.18$.
Black-Litterman recovers implied $\Pi$ in $[5.64\%, 10.49\%]$ vs sample
$[12.64\%, 30.90\%]$. ERC achieves $\mathrm{HHI}_{\text{risk}} = 0.067
\approx 1/N$ by construction against $0.069$ for $1/N$ and $0.318$ for
Markowitz. \textbf{Contribution.} The paper unifies five models under
a single estimation pipeline with comparable per-asset weights. The
implementation is open-source, the data are public, the platform is
free.
\vspace{0.4cm}\noindent\textbf{Keywords:}
Markowitz; DeMiguel-Garlappi-Uppal; Kahneman; Black-Litterman; risk
parity; Ledoit-Wolf; Jorion; tangency portfolio; Medallion lakehouse;
Databricks; B3; IBOV; reproducibility.
\end{singlespace}

\newpage

% ---------- SUMÁRIO ------------------------------------------------------
\tableofcontents
\newpage

% =====================================================================
\section{Introdução}\label{sec:intro}

A teoria moderna de carteiras estabeleceu-se em meados do século XX
com o trabalho seminal de \textbf{Harry Markowitz} (1952), que mostrou
\textit{como} um investidor racional deveria alocar riqueza entre $n$
ativos arriscados sob a hipótese de que utilidade depende apenas da
média e da variância dos retornos. A inovação técnica foi modesta
\textemdash{} um problema de otimização quadrática com restrições
lineares \textemdash{} mas a inovação conceitual foi transformadora:
pela primeira vez o risco da carteira era tratado como uma propriedade
\textit{emergente} da matriz de covariância $\Sigma$, não como uma
soma simples dos riscos individuais. Eficiência passou a significar
\textit{trade-off ótimo} entre retorno esperado e risco, e o ferramental
gráfico da fronteira de média-variância passou a ser o vocabulário
universal de gestores.

O modelo de Markowitz tornou-se ortodoxia acadêmica e operacional.
Foi estendido por \textbf{Tobin} (1958) com a introdução de um ativo
livre de risco, gerando a Linha do Mercado de Capitais (CML);
generalizado por \textbf{Sharpe} (1964) e Lintner (1965) ao CAPM, que
abstrai a estrutura de covariância para uma única dimensão de risco
sistemático ($\beta$); refinado por \textbf{Merton} (1972), que
forneceu a forma fechada da fronteira eficiente; estendido a tempo
contínuo (Merton, 1973), a múltiplos fatores (Fama-French, 1993; Carhart,
1997), e a horizontes de longo prazo (Merton, 1969). \textbf{Markowitz},
\textbf{Sharpe} e \textbf{Miller} compartilharam o prêmio Nobel de
Ciências Econômicas de 1990 pelo desenvolvimento da teoria moderna
de finanças.

Paralelamente à expansão teórica, a aplicação prática da otimização
Markowitz revelou-se sistematicamente decepcionante. Já em
\textbf{Michaud} (1989), \textit{The Markowitz Optimization Enigma is
Not Optimized Optimization}, documenta-se que carteiras ótimas
construídas a partir de estimadores amostrais $\hat{\mu}, \hat{\Sigma}$
são extremamente sensíveis a perturbações pequenas em $\hat{\mu}$,
resultando em soluções concentradas em poucos ativos, com pesos
instáveis no tempo, e desempenho \textit{out-of-sample} frequentemente
inferior a regras ingênuas. Michaud cunha o termo \textit{error
maximization}: o otimizador \emph{amplifica} os erros de estimativa em
$\hat{\mu}$, porque a função-objetivo $w^\top \mu - \tfrac{\delta}{2}
w^\top \Sigma w$ é não-linear em $\mu$ e premia desvios extremos.

Duas décadas depois, \textbf{DeMiguel, Garlappi e Uppal} (2009),
\textit{Optimal Versus Naive Diversification: How Inefficient Is the
1/N Portfolio Strategy?}, formalizam empiricamente a crítica em 14
\textit{datasets} clássicos (FF industries, MSCI countries, S\&P
sectors) e mostram que \textit{nenhuma} entre 14 variantes testadas de
otimização (Markowitz amostral, mínima-variância, Bayes-Stein,
Black-Litterman, \textit{kernel-shrinkage}) supera robustamente o
$1/N$ ingênuo após custos de transação \textit{e} em horizonte longo.
A interpretação formal é direta: o erro padrão da média amostral é
$\sigma/\sqrt{T}$; para vol típica de ações brasileiras
$\sim 30\%$ a.a.\ e janela de 5 anos ($T = 1260$), o desvio em
$\hat{\mu}$ é da ordem de $\pm 8{,}5$ pontos percentuais, comparável
ao próprio prêmio de risco do mercado.

Em paralelo, \textbf{Daniel Kahneman} \textemdash{} prêmio Nobel de
Ciências Econômicas de 2002 \textemdash{} desenvolve, em colaboração
com Amos Tversky, uma crítica \textit{comportamental} ao paradigma da
gestão ativa. Em \textit{Thinking, Fast and Slow} (2011), Kahneman
documenta um estudo realizado em 1984 com uma firma de \textit{wealth
management} que possuía 25 anos de avaliações individuais para seus 28
consultores. A correlação dos rankings de retorno entre pares de anos
consecutivos \textemdash{} 378 pares no total \textemdash{} foi
\textbf{0{,}01} em média. Kahneman batiza o achado de \textit{The
Illusion of Skill}: a indústria financeira é remunerada por uma
persistência de \textit{performance} que, sob qualquer teste estatístico
razoável, não se distingue de sorte.

A literatura responde a Michaud e DGU com três famílias de mitigações.
A primeira, \textit{shrinkage estatístico}: \textbf{Jorion} (1986)
propõe encolhimento Bayes-Stein de $\hat{\mu}$ na direção da
média comum; \textbf{Ledoit e Wolf} (2003, 2004) propõem encolhimento
data-driven de $\hat{\Sigma}$ na direção de alvos paramétricos
(identidade, correlação constante). A segunda, \textit{regularização
bayesiana}: \textbf{Black e Litterman} (1992) substituem a média
amostral pela média implícita pelo equilíbrio CAPM
$\Pi = \delta \Sigma w_{\text{mkt}}$ e combinam essa prior com
\textit{views} explicitamente declaradas, com matriz de incerteza
$\Omega$ calibrável. \textbf{He e Litterman} (1999) e
\textbf{Idzorek} (2005) refinam a calibração de $\Omega$. A terceira,
\textit{descarte de $\mu$}: \textbf{Maillard, Roncalli e Teïletche}
(2010) propõem o portfólio de Paridade de Risco (\textit{Equal Risk
Contribution}), em que cada ativo contribui igualmente para a
variância total, dispensando inteiramente o estimador de retorno
esperado e usando apenas $\Sigma$.

\subsection*{Objetivo e contribuição deste trabalho}

Este \textit{working paper} replica os cinco modelos
\textemdash{} Markowitz puro, $1/N$ de DGU, ilusão de skill de
Kahneman, Black-Litterman e Paridade de Risco \textemdash{} sobre o
universo brasileiro de ações listadas na B3, em janela
$2000$\textendash$2026$, sob um único pipeline integrado de estimação.
Os cinco modelos são contrastados ativo-a-ativo: os mesmos pesos
ranqueados, sob a mesma matriz $\Sigma$ Ledoit-Wolf, com a mesma taxa
livre de risco $r_f$ e o mesmo prêmio de risco $\mathrm{ERP}$ ancorando
o macro-shrinkage. A reprodução é integral: o código é
\textit{open-source} (TypeScript no \textit{front-end}, Python no
pipeline Databricks), os dados são públicos (Yahoo Finance + Banco
Central do Brasil), e a plataforma é gratuita (Databricks Free
Edition).

\subsection*{Estrutura do artigo}

A Seção~\ref{sec:teoria} consolida o referencial teórico, citando
os trabalhos primários e os \textit{prêmios Nobel} relacionados.
A Seção~\ref{sec:dados} descreve o conjunto de dados:
o universo curado de 982 \textit{tickers}, o processo de ingestão via
Yahoo Finance, as 38 constituintes do IBOV efetivamente disponíveis,
a janela temporal e o tratamento de retornos. A Seção~\ref{sec:metodo}
formaliza o pipeline de estimação: Ledoit-Wolf,
Jensen-correction, Jorion shrinkage, macro-anchor, \textit{long-only}
\textit{active-set}. A Seção~\ref{sec:resultados}
apresenta os resultados empíricos das cinco abordagens, com tabelas e
estatísticas-resumo. A Seção~\ref{sec:discussao} discute as
implicações, em particular a observação \textit{contra-intuitiva} de
que Markowitz superou o $1/N$ na janela $2024$\textendash$2026$
\textemdash{} e por que isso não invalida a tese DGU.
A Seção~\ref{sec:conclusao} fecha com considerações finais
e direções de pesquisa futura.

% =====================================================================
\section{Referencial teórico}\label{sec:teoria}

\subsection{O problema de Markowitz (1952)}

A formulação canônica do problema de Markowitz, em sua forma quadrática
sem restrição de ativo livre de risco, é
\begin{equation}\label{eq:markowitz-min-var}
\min_{w \in \mathbb{R}^n} \quad \tfrac{1}{2} w^\top \Sigma w
\quad \text{sujeito a} \quad w^\top \mu = \mu^*, \quad w^\top \mathbf{1} = 1,
\end{equation}
onde $w \in \mathbb{R}^n$ é o vetor de pesos, $\mu \in \mathbb{R}^n$ o
vetor de retornos esperados, $\Sigma \in \mathbb{R}^{n \times n}$ a
matriz de covariância (positivo-definida), $\mu^*$ o retorno-alvo e
$\mathbf{1}$ o vetor de uns. A solução analítica de Merton (1972) define
três quantidades-base $A = \mathbf{1}^\top \Sigma^{-1} \mathbf{1}$,
$B = \mathbf{1}^\top \Sigma^{-1} \mu$, $C = \mu^\top \Sigma^{-1} \mu$,
$D = AC - B^2$, com as quais a fronteira eficiente fica caracterizada
por
\begin{equation}
\sigma^2(\mu^*) = \frac{A (\mu^*)^2 - 2 B \mu^* + C}{D},
\end{equation}
e a carteira de mínima variância tem $\mu_{\text{mv}} = B/A$,
$\sigma^2_{\text{mv}} = 1/A$.

Quando se introduz um ativo livre de risco com taxa $r_f$, a fronteira
eficiente colapsa para uma reta na seguinte sentido: \textit{toda}
carteira eficiente é uma combinação linear do ativo livre de risco e
da \textbf{carteira de tangência}
\begin{equation}\label{eq:tangency}
w_T = \frac{\Sigma^{-1} (\mu - r_f \mathbf{1})}{\mathbf{1}^\top \Sigma^{-1} (\mu - r_f \mathbf{1})}.
\end{equation}
A inclinação da reta tangente é o $\mathrm{Sharpe} = (\mu_T - r_f) / \sigma_T$
da carteira de tangência \textemdash{} o preço de mercado do risco.

A solução é \textit{conceitualmente} elegante mas \textit{empiricamente}
frágil, pelas razões já antecipadas: $\hat{\mu}$ é um estimador ruidoso
e $\Sigma^{-1}$ amplifica esse ruído quando $\Sigma$ é mal-condicionada
(condição agravada por $N \gtrsim T$).

\subsection{CAPM e a estrutura de equilíbrio (Sharpe, 1964)}

\textbf{William Sharpe} (1964) \textemdash{} prêmio Nobel 1990
\textemdash{} agregou as preferências individuais de Markowitz num
modelo de equilíbrio. Sob hipóteses fortes (preferências quadráticas
homogêneas, expectativas homogêneas, mercados sem fricção, taxa $r_f$
única), todos os investidores optam pela mesma carteira de
tangência. Em equilíbrio, essa carteira coincide com a
\textbf{carteira de mercado}: $w_T = w_{\text{mkt}}$. O preço de
qualquer ativo $i$ é determinado pela sua contribuição ao risco
sistemático,
\begin{equation}\label{eq:capm}
E[r_i] - r_f = \beta_i (E[r_{\text{mkt}}] - r_f),
\quad \beta_i = \frac{\mathrm{Cov}(r_i, r_{\text{mkt}})}{\mathrm{Var}(r_{\text{mkt}})}.
\end{equation}

A relevância do CAPM para este trabalho é dupla: (a) ele fornece a
base teórica para o modelo Black-Litterman, que inverte \eqref{eq:capm}
para recuperar os retornos implícitos $\Pi = \delta \Sigma w_{\text{mkt}}$
a partir da carteira de mercado observada; e (b) ele estabelece a
carteira de mercado como o benchmark \textit{natural} para qualquer
avaliação de gestão ativa \textemdash{} dispensar o índice por
otimização sem informação privada deveria ser, em equilíbrio, neutro
ou prejudicial.

\subsection{A crítica de Michaud (1989)}

\textbf{Robert Michaud} (1989), em \textit{The Markowitz Optimization
Enigma is Not Optimized Optimization}, publica na \textit{Financial
Analysts Journal} o argumento que se tornaria referência para todo o
debate posterior sobre fragilidade do MV. O artigo demonstra
empiricamente que perturbações pequenas em $\hat{\mu}$ \textemdash{} da
ordem de 25 pontos-base, dentro do erro padrão usual da média
amostral \textemdash{} podem inverter completamente a composição da
carteira ótima, alternando pesos entre $+100\%$ e $-100\%$ em ativos
individuais. Michaud chama esse comportamento de \textit{error
maximization}: a função-objetivo de \eqref{eq:markowitz-min-var}
premia justamente os ativos cujo $\hat{\mu}_i$ está mais inflado por
sorte amostral, porque a otimização não tem como distinguir sinal de
ruído quando a janela é curta. A solução é não usar a otimização puro
forma; alternativas plausíveis incluem \textit{resampling} (Michaud,
1998), \textit{shrinkage}, \textit{Bayesian portfolio choice} e
restrições de variabilidade.

A formalização posterior do achado de Michaud foi feita por
\textbf{Best e Grauer} (1991), também na FAJ, que provam analiticamente
que a sensibilidade dos pesos ótimos a perturbações em $\mu$ é
proporcional a $\Sigma^{-1}$. Para covariâncias com autovalores
pequenos (ativos quase-redundantes), $\Sigma^{-1}$ explode em norma e
a sensibilidade dos pesos diverge.

\subsection{Shrinkage estatístico: Jorion (1986) e Ledoit-Wolf (2003, 2004)}

\textbf{Philippe Jorion} (1986), em \textit{Bayes-Stein Estimation
for Portfolio Analysis}, propõe encolher $\hat{\mu}$ na direção da
média ponderada de todos os ativos, $\bar{\mu}$, com intensidade
$\psi^* \in (0, 1)$ calibrada para minimizar o risco quadrático:
\begin{equation}\label{eq:jorion}
\hat{\mu}^{\text{JS}} = (1 - \psi^*) \hat{\mu} + \psi^* \bar{\mu}
\cdot \mathbf{1},
\quad \psi^* = \frac{(N+2)}{(N+2) + T \cdot (\hat{\mu} - \bar{\mu})^\top \Sigma^{-1} (\hat{\mu} - \bar{\mu})}.
\end{equation}
Em horizontes curtos ($T$ pequeno), $\psi^* \to 1$ e o estimador é
toda média comum; em horizontes longos ($T$ grande), $\psi^* \to 0$ e
recuperamos $\hat{\mu}$.

\textbf{Olivier Ledoit e Michael Wolf} (2003, 2004) generalizam o
shrinkage para a matriz de covariância, propondo:
\begin{equation}\label{eq:ledoit-wolf}
\hat{\Sigma}^{\text{LW}} = (1 - \delta^*) \hat{\Sigma} + \delta^* F,
\end{equation}
onde $F$ é um alvo paramétrico (identidade escalada, correlação
constante, ou modelo de um fator) e $\delta^* \in (0, 1)$ é o nível de
encolhimento ótimo, derivado em forma fechada como minimização do erro
quadrático esperado entre $\hat{\Sigma}^{\text{LW}}$ e $\Sigma$. O
estimador resultante é sempre positivo-definido (mesmo quando
$\hat{\Sigma}$ não é, no regime $N > T$) e domina o sample em risco
quadrático.

A combinação Jorion (para $\mu$) + Ledoit-Wolf (para $\Sigma$) é o
\textit{stack} padrão de \textit{shrinkage} em finanças quantitativas
contemporâneas. É o adotado neste trabalho, com extensão adicional do
\textit{macro-anchor} descrito na Seção~\ref{sec:metodo}.

\subsection{O argumento empírico: DeMiguel, Garlappi e Uppal (2009)}

\textbf{Victor DeMiguel, Lorenzo Garlappi e Raman Uppal} (2009)
publicam na \textit{Review of Financial Studies} a peça que reorganiza
o debate. O artigo testa 14 modelos de otimização \textemdash{}
Markowitz amostral, mínima variância, Bayes-Stein, Black-Litterman
com vários valores de $\tau$, MacKinlay-Pástor, encolhimento da
covariância, encolhimento da média, modelo de fatores, abordagens
\textit{constrained} \textemdash{} contra o $1/N$ ingênuo em 14
\textit{datasets} clássicos. A métrica é Sharpe \textit{out-of-sample},
após custos de transação, em janela rolante.

O resultado é taxativo: \textit{nenhum} dos 14 modelos supera robustamente
o $1/N$. Em \textit{datasets} maiores (mais ativos por janela fixa), o
ruído domina o sinal e o $1/N$ vence consistentemente. Em
\textit{datasets} menores e janelas longas, alguns modelos empatam mas
com volatilidade alta. A interpretação econômica é direta: o
$\hat{\mu}$ amostral é estatisticamente ruim demais para que a
otimização agregue valor após o custo de implementação
(\textit{turnover}, \textit{tracking error}, complexidade operacional).

DGU oferece também um cálculo de tamanho de amostra: para que a
otimização de média-variância supere o $1/N$ com 95\% de confiança,
seria necessário $T \approx 3.000$ anos de história ininterrupta com
$N = 25$ ativos. É um cálculo prático: a vida humana não é longa o
suficiente para que $\hat{\mu}$ amostral seja confiável.

\subsection{A crítica comportamental: Kahneman (1984/2011)}

\textbf{Daniel Kahneman} \textemdash{} prêmio Nobel de Ciências
Econômicas de 2002 \textemdash{}, em colaboração com Amos Tversky,
desenvolve a teoria do prospecto (Kahneman \& Tversky, 1979) e, num
programa de pesquisa de várias décadas, documenta sistematicamente os
\textit{vieses cognitivos} que afetam tomada de decisão sob incerteza.
A síntese final aparece em \textit{Thinking, Fast and Slow} (2011),
livro de divulgação científica e referência canônica.

Para finanças, a contribuição chave é o capítulo \textit{The Illusion
of Stock-Picking Skill}, em que Kahneman relata um estudo que conduziu
em 1984. Uma firma de \textit{wealth management} (não-identificada)
forneceu 25 anos de avaliações individuais para seus 28 consultores,
com bônus calculados por ranking de \textit{performance} relativa.
Kahneman calculou todas as correlações de ranking entre pares de anos
consecutivos \textemdash{} $\binom{25-1}{1} = 24$ correlações
\textit{por consultor}, totalizando $28 \cdot 24 / 2 = 336$ comparações
de pares de anos. A média foi $0{,}01$. Persistência de
\textit{performance} ano-a-ano: estatisticamente nula.

A interpretação de Kahneman é o que ele chama de \textit{illusion of
validity}: o setor financeiro é \textbf{pago para entregar uma habilidade
que estatisticamente não existe}. O argumento se conecta a DGU mas é
distinto: DGU é uma crítica \textit{matemática} ao otimizador, Kahneman
é uma crítica \textit{comportamental} ao mercado de gestão ativa.
Convergem na mesma conclusão pragmática: para o pequeno investidor,
$1/N$ ou índice de mercado supera tentar ser esperto.

A formalização estatística do vies de Kahneman foi feita por
\textbf{Kan e Smith} (2008), \textit{The Distribution of the Sample
Minimum-Variance Frontier}, na \textit{Management Science}. Os autores
mostram que o Sharpe \textit{in-sample} de um estimador de carteira de
tangência é viciado para cima por um fator
\begin{equation}\label{eq:kan-smith}
\mathbb{E}[\widehat{\mathrm{Sharpe}}^2_{\text{in}}] \;-\;
\mathrm{Sharpe}^2_{\text{true}} \;\approx\; \frac{N}{T}
\;+\; O\!\left(\frac{N^2}{T^2}\right),
\end{equation}
para $N$ ativos e janela $T$. Em $N = 30, T = 504$ (dois anos),
o vies é da ordem de $30/504 = 0{,}06$ em $\mathrm{Sharpe}^2$, o que
desloca um Sharpe verdadeiro de $0{,}5$ para um Sharpe declarado de
$\sqrt{0{,}25 + 0{,}06} \approx 0{,}56$ \textemdash{} ou seja,
$+0{,}06$ em pontos de Sharpe linear, um vies sistemático que cresce
com $N$ e cai com $T$.

\subsection{O remédio bayesiano: Black-Litterman (1992)}

\textbf{Fischer Black e Robert Litterman} (1992), em
\textit{Global Portfolio Optimization} (\textit{Financial Analysts
Journal}), trabalho desenvolvido na Goldman Sachs Fixed Income
Research, propõem o que se tornaria o \textit{framework} bayesiano
canônico de seleção de carteira. A construção tem duas etapas.

Primeiro, em vez de partir de $\hat{\mu}$ amostral, os autores partem
dos \textbf{retornos implícitos pelo equilíbrio CAPM}:
\begin{equation}\label{eq:bl-pi}
\Pi = \delta \, \Sigma \, w_{\text{mkt}},
\end{equation}
onde $w_{\text{mkt}}$ é a carteira de mercado observada e
$\delta$ é o coeficiente de aversão a risco do investidor
representativo (tipicamente $\delta \in [2, 3]$; He-Litterman usam
$\delta = 2{,}5$). $\Pi$ é o vetor de retornos que tornaria
$w_{\text{mkt}}$ a solução ótima de média-variância
\eqref{eq:tangency} \textemdash{} é uma \textbf{inversão} do CAPM.

Segundo, os autores tratam $\Pi$ como a média da \textit{prior}
bayesiana de $\mu$:
\begin{equation}\label{eq:bl-prior}
\mu \sim \mathcal{N}(\Pi, \, \tau \Sigma),
\end{equation}
e permitem que o investidor adicione \textit{views} lineares
$P \mu = Q + \varepsilon$, $\varepsilon \sim \mathcal{N}(\mathbf{0},
\Omega)$, em que cada linha de $P \in \mathbb{R}^{K \times N}$ é a
combinação linear de ativos sobre a qual a \textit{view} incide, $Q_k$
é o retorno esperado predito, e $\Omega$ é diagonal com $\Omega_{kk}$
inversamente proporcional à confiança do investidor na \textit{view}
$k$. A combinação bayesiana resultante é
\begin{equation}\label{eq:bl-posterior}
\mu_{\text{BL}} = \bigl[(\tau \Sigma)^{-1} + P^\top \Omega^{-1} P\bigr]^{-1}
\bigl[(\tau \Sigma)^{-1} \Pi + P^\top \Omega^{-1} Q\bigr],
\end{equation}
com matriz de covariância da estimativa
\begin{equation}\label{eq:bl-cov}
M = \bigl[(\tau \Sigma)^{-1} + P^\top \Omega^{-1} P\bigr]^{-1},
\end{equation}
e covariância \textit{efetiva} usada no otimizador
$\Sigma_{\text{BL}} = \Sigma + M$.

Sem \textit{views}, $\mu_{\text{BL}} = \Pi$ exatamente; a tangência sob
$(\mu_{\text{BL}}, \Sigma_{\text{BL}})$ recupera $w_{\text{mkt}}$
(modulo o termo $\tau$). É o ponto teórico mais importante do modelo:
\textbf{sem opinião, o investidor deveria manter o índice}. Com
\textit{views} de alta confiança, o posterior tilta na direção
prescrita. Com \textit{views} de baixa confiança, o posterior fica
próximo de $\Pi$.

\textbf{He e Litterman} (1999), \textit{The Intuition Behind
Black-Litterman Model Portfolios}, refinam a calibração de $\Omega$ e
fornecem uma cartilha aplicada. \textbf{Idzorek} (2005),
\textit{A Step-by-Step Guide to the Black-Litterman Model}, propõe
calibrar $\Omega$ a partir de uma confiança percentual explicitamente
declarada pelo gestor, mapeando $0\%$ confiança $\to \Omega \to \infty$
e $100\%$ confiança $\to \Omega \to 0$.

\subsection{Jogando fora $\mu$: Paridade de Risco (Maillard et al., 2010)}

\textbf{Sébastien Maillard, Thierry Roncalli e Jérôme Teïletche} (2010),
\textit{The Properties of Equally Weighted Risk Contribution Portfolios}
(\textit{Journal of Portfolio Management}), oferecem uma resposta
radical ao problema de estimação de $\mu$: \textbf{descartá-lo
inteiramente}. O argumento é que, dada a evidência de DGU sobre a
fragilidade de $\hat{\mu}$, vale mais focar em uma quantidade que
\textit{pode} ser estimada com robustez \textemdash{} a covariância
$\Sigma$ \textemdash{} e definir o portfólio em termos dela apenas.

A contribuição de risco do ativo $i$ é definida como a fração da
variância total da carteira atribuível a $i$:
\begin{equation}\label{eq:rc}
\mathrm{RC}_i(w) \;=\; \frac{w_i (\Sigma w)_i}{w^\top \Sigma w},
\quad \sum_{i=1}^N \mathrm{RC}_i = 1.
\end{equation}
A derivada parcial de $\sigma_p = \sqrt{w^\top \Sigma w}$ em relação a
$w_i$ é $(\Sigma w)_i / \sigma_p$; portanto $w_i \cdot \partial
\sigma_p / \partial w_i$ é a contribuição \textit{euleriana} de $i$
para a volatilidade total. Dividindo por $\sigma_p$ obtém-se
\eqref{eq:rc} como contribuição fracional para a variância.

A carteira de \textbf{Equal Risk Contribution (ERC)} é o vetor de
pesos que satisfaz $\mathrm{RC}_i = 1/N$ para todo $i$. Maillard et
al.\ provam que (a) a solução existe e é única para $\Sigma$
positivo-definida, (b) sua volatilidade fica entre a do $1/N$ e a do
mínimo-variância, e (c) é o único minimizador da função convexa
\begin{equation}\label{eq:erc-objective}
\min_{w \geq 0} \quad \tfrac{1}{2} w^\top \Sigma w
\;-\; \frac{1}{N} \sum_{i=1}^N \ln(w_i),
\end{equation}
sujeito apenas à positividade de $w$. O termo $-\ln(w_i)$ é uma
\textbf{barreira logarítmica} que impede $w_i \to 0$ (porque $\ln 0 =
-\infty$). O gradiente da função-objetivo é $\Sigma w - (1/N) (1/w_i)$,
e a condição de primeira ordem é $w_i (\Sigma w)_i = $ constante,
equivalente a $\mathrm{RC}_i = 1/N$.

O caso especial em que $\Sigma$ é diagonal admite solução fechada:
\begin{equation}\label{eq:inv-vol}
w_i^{\text{inv-vol}} \;\propto\; \frac{1}{\sigma_i},
\quad \sigma_i = \sqrt{\Sigma_{ii}},
\end{equation}
que é a \textit{paridade de volatilidade} ou \textit{naive risk
parity}, implementada pela maioria dos produtos comerciais.

O paradigma de Paridade de Risco ganhou popularidade fora da academia
com o \textbf{All-Weather Portfolio} da Bridgewater Associates, lançado
por Ray Dalio em 1996. O All-Weather aloca igualmente em quatro
``quadrantes'' macroeconômicos (crescimento alto/baixo $\times$
inflação alta/baixa), calibrando pesos para que cada quadrante
contribua igualmente para a vol total. \textbf{Roncalli} (2013),
\textit{Introduction to Risk Parity and Budgeting} (Chapman \&
Hall/CRC), documenta a formalização acadêmica posterior e as
extensões a múltiplos níveis (Hierarchical Risk Parity, López de
Prado 2016).

\subsection{Síntese do referencial}

Os cinco modelos contrastados neste artigo se distribuem em três
níveis epistêmicos:
\begin{enumerate}
\item \textbf{Markowitz puro} adota $\hat{\mu}, \hat{\Sigma}$ amostrais
em \eqref{eq:tangency} sem nenhuma proteção contra ruído. Serve como
\textit{benchmark teórico}.
\item \textbf{$1/N$ ingênuo} (DGU) ignora $\Sigma$ além da contagem $N$
e ignora $\mu$ inteiramente. Serve como \textit{benchmark empírico}.
\item Os três restantes \textemdash{} \textbf{Kahneman} (experimento de
ilusão de skill), \textbf{Black-Litterman}, e \textbf{Paridade de
Risco} \textemdash{} oferecem três respostas distintas ao
problema de estimação: respectivamente, \textit{quantificar} a
diferença entre promessa e entrega, \textit{substituir} a média
amostral pela média de equilíbrio, e \textit{descartar} a média.
\end{enumerate}

Este artigo replica todos eles sob um pipeline unificado e exibe os
resultados ativo-a-ativo. A próxima Seção descreve os dados.

% =====================================================================
\section{Dados e materiais}\label{sec:dados}

\subsection{A plataforma \textit{lakehouse}}

Os dados que sustentam este artigo são produzidos por uma plataforma
\textit{lakehouse} Medallion \textit{open-source}, hospedada no
Databricks Free Edition do autor. A arquitetura é a canônica
\textbf{bronze} $\rightarrow$ \textbf{silver} $\rightarrow$
\textbf{gold} (Armbrust et al., 2021; Databricks, 2024), com três
camadas:

\textbf{Bronze}: ingestão bruta, sem transformação. Inclui (i)
\texttt{bronze.b3\_universe} (universo curado de ativos B3, carregado
a partir de uma planilha CSV versionada), (ii) \texttt{bronze.b3\_ohlcv\_raw}
(OHLCV diário ajustado por dividendos e splits, obtido via
Yahoo Finance), e (iii) \texttt{bronze.b3\_index\_members} (composição
do IBOV ao longo do tempo).

\textbf{Silver}: limpeza, deduplicação e enriquecimento. Inclui
\texttt{silver.b3\_ohlcv\_adjusted} (OHLCV final ajustado),
\texttt{silver.b3\_ticker\_dim} (dimensão de tickers com setor B3,
\textit{listed\_from}/\textit{listed\_to}, prior tickers para fusões
e \textit{renamings}), e \texttt{silver.b3\_index\_members\_long}
(formato \textit{slowly-changing dimension} para a composição IBOV).

\textbf{Gold}: agregados de consumo. Para este artigo, os relevantes
são \texttt{gold.returns\_wide} (matriz $T \times N$ de log-retornos
diários em formato \textit{wide} \textemdash{} uma coluna por
\textit{ticker}, uma linha por data de pregão),
\texttt{gold.cov\_matrix\_5y} (matriz $\Sigma$ Ledoit-Wolf anualizada
sobre os últimos 1260 dias úteis, formato \textit{long}),
\texttt{gold.ibov\_overview} (composição IBOV \textit{as-of} corrente
com pesos, retorno YTD e contribuição), e
\texttt{gold.kpis\_per\_ticker} (KPIs individuais por
\textit{ticker}: $\sigma$ anualizada, drawdown máximo, Sharpe vs CDI,
$N$ observações).

Os artefatos JSON resultantes da camada gold são exportados para um
\textit{volume} no Databricks (\texttt{/Volumes/finance\_prd/gold/artifacts/})
e replicados ao GitHub Pages a cada \textit{deploy} via workflow CI
(\texttt{.github/workflows/deploy-pages.yml}). A \textit{front-end} em
Next.js carrega esses JSON estáticos no \textit{build time} e os usa
para todas as visualizações.

\subsection{O universo curado: 982 \textit{tickers}}

O universo curado de ativos é mantido em
\texttt{data/ticker\_universe.csv}, versionado junto ao código do
projeto. A última versão contém \textbf{982 \textit{tickers}}, cobrindo
ações ordinárias (ON), preferenciais (PN), units, BDRs, ETFs e fundos
imobiliários listados na B3 nos últimos 25 anos. A planilha inclui
metadata curatorial: setor B3 e subsetor, \texttt{listed\_from}
(data de listagem original), \texttt{listed\_to} (NULL se ainda ativo),
\texttt{prior\_tickers} (\textit{array} de tickers anteriores, para
ativos que passaram por \textit{rename} ou fusão \textemdash{} ex.\
BRFS3 = SADIA3 $|$ PRGA3; VBBR3 = BRDT3; RAIL3 = ALLL3), e CNPJ
quando conhecido.

\subsection{A ingestão Yahoo Finance: 468 \textit{tickers} efetivos}

A ingestão OHLCV diária via Yahoo Finance retorna dados para apenas
\textbf{468 dos 982 \textit{tickers}} curados. Os \textbf{514 ausentes}
distribuem-se em três grupos: (i) ativos delistados há mais de 10 anos
cuja \textit{API} de Yahoo retorna vazio com o \textit{ticker} atual;
(ii) ativos que passaram por \textit{rename} ou fusão e cuja história
está sob o nome antigo no Yahoo \textemdash{} a ingestão recupera o
nome novo, encontra histórico recente curto, e classifica como
``\textit{insufficient history}''; e (iii) ativos pouco-líquidos cujo
\textit{ticker} Yahoo é estruturalmente diferente do código B3 (BDRs,
algumas \textit{units}).

A consequência prática é que o universo \textbf{efetivamente
investível} para este trabalho é de 468 \textit{tickers}. A janela
2000-01-03 a 2026-05-25 oferece 6.791 dias úteis de história, mas
muitos \textit{tickers} têm histórico parcial (IPOs entre 2010 e 2024).

\subsection{O IBOV e as 38 constituintes disponíveis}

A composição corrente do IBOV (\textit{as-of} 2026-05-25) tem
\textbf{51 ativos} ativos em \texttt{silver.b3\_index\_members\_long}.
A camada gold (\texttt{gold.ibov\_overview}), porém, agrega apenas
\textit{tickers} com preço disponível em \texttt{silver.b3\_ohlcv\_adjusted}.
Treze constituintes correntes do IBOV \textemdash{} BRFS3, CCRO3,
CPLE6, ELET3, ELET6, EMBR3, JBSS3, NTCO3, RAIL3, SUZB3, TIMS3, VBBR3,
VIVT3 \textemdash{} não retornam dados de Yahoo com o \textit{ticker}
atual (provavelmente pela mesma raiz dos 514 ausentes do universo
amplo). A composição IBOV efetivamente disponível para uso na
otimização Markowitz e Black-Litterman é, portanto, de \textbf{38
constituintes}, cobrindo aproximadamente 81\% do peso real do índice.

A tabela~\ref{tab:ibov-top10} apresenta as dez maiores ponderações.

\begin{table}[H]
\centering
\caption{IBOV: dez maiores constituintes disponíveis no gold (peso, setor B3).}\label{tab:ibov-top10}
\begin{tabular}{lrl}
\toprule
\textbf{Ticker} & \textbf{Peso} & \textbf{Setor B3} \\
\midrule
VALE3   & 9{,}08\% & Materiais Básicos \\
PETR4   & 6{,}73\% & Petróleo Gás e Biocombustíveis \\
ITUB4   & 6{,}70\% & Financeiro \\
PETR3   & 3{,}87\% & Petróleo Gás e Biocombustíveis \\
BBAS3   & 3{,}53\% & Financeiro \\
B3SA3   & 3{,}19\% & Financeiro \\
WEGE3   & 2{,}97\% & Bens Industriais \\
ABEV3   & 2{,}88\% & Consumo Não Cíclico \\
BBDC4   & 2{,}76\% & Financeiro \\
BPAC11  & 2{,}45\% & Financeiro \\
\bottomrule
\end{tabular}
\end{table}

A concentração setorial é notável: dos dez maiores, \textbf{cinco são
financeiros} (ITUB4, BBAS3, B3SA3, BBDC4, BPAC11) somando $18{,}63\%$
do peso, e \textbf{dois são petróleo} (PETR4, PETR3) somando
$10{,}60\%$. Materiais Básicos (VALE3) sozinho é $9{,}08\%$.

\subsection{Taxa livre de risco}

A taxa livre de risco usada em todo o pipeline é a média histórica do
CDI, obtida via Banco Central do Brasil (BCB SGS série 12). No
último \textit{snapshot} disponível, o valor é
$r_f = 12{,}90\% \text{ a.a.}$, cobrindo 27 anos de história
(1999\textendash 2026).

\subsection{Período do estudo}

A janela de trabalho é \textbf{2000-01-03 a 2026-05-25}, totalizando
6.791 datas de pregão B3. As estimações com janela móvel restringem-se
às datas em que o \textit{ticker} avaliado tem cobertura completa
($\geq 95\%$ de dias não-nulos para janelas curtas, ou cobertura
adaptativa $\geq \max\{0{,}60, \; 0{,}95 \cdot \max_i \mathrm{cobertura}_i\}$
para janelas longas \textemdash{} ver seção~\ref{sec:metodo}).

% =====================================================================
\section{Metodologia}\label{sec:metodo}

\subsection{Estimação da matriz de covariância}\label{sec:metodo:sigma}

A matriz de covariância $\hat{\Sigma}$ é estimada via \textbf{Ledoit-Wolf
(2004)} com alvo de correlação constante. Em forma compacta:
\begin{equation}
\hat{\Sigma}^{\text{LW}} = (1 - \delta^*) S + \delta^* F,
\end{equation}
onde $S = \tfrac{1}{T-1} \sum_t (r_t - \bar{r})(r_t - \bar{r})^\top$ é
a covariância amostral, $F$ é o alvo construído a partir de
$\bar{\rho} \sqrt{\Sigma_{ii} \Sigma_{jj}}$ para $i \neq j$ (com
$\bar{\rho}$ = média das correlações amostrais fora da diagonal), e
$\delta^* \in [0, 1]$ é a intensidade de \textit{shrinkage} ótima
calculada em forma fechada como
$\delta^* = \max\{0, \min\{1, \pi/(\gamma T)\}\}$, sendo $\pi$ a soma
das variâncias assintóticas dos elementos de $S$ e $\gamma$ a soma dos
quadrados dos desvios $(F - S)^2$. A implementação está em
\texttt{pipelines/notebooks/gold/cov\_matrix.py}.

Antes de Ledoit-Wolf, a matriz é filtrada para incluir apenas
\textit{tickers} com cobertura suficiente na janela. A regra é
\textit{hybrid}: \textbf{para janelas curtas} (1y, 5y, 10y), exigimos
cobertura absoluta $\geq 95\%$; \textbf{para janelas longas} (15y, 20y,
\textit{full}), a regra é adaptativa,
\begin{equation}\label{eq:cov-threshold}
\text{threshold} = \max\bigl\{ 0{,}60, \;\; 0{,}95 \cdot \max_c \text{cobertura}_c \bigr\},
\end{equation}
para acomodar o fato de que mesmo o \textit{ticker} mais antigo no
universo expandido pode ter cobertura inferior a $95\%$ em 15 ou 20
anos (porque IPOs recentes introduzem datas que ele não cobre). Após
filtragem, dropamos linhas com qualquer NaN entre os \textit{tickers}
selecionados para garantir painel retangular limpo para o cálculo de
$S$. A tabela~\ref{tab:cov-windows} resume os \textit{n\_valid} por
janela.

\begin{table}[H]
\centering
\caption{Estatísticas de cobertura por janela de covariância (universo 468).}\label{tab:cov-windows}
\begin{tabular}{lrrr}
\toprule
\textbf{Janela} & \textbf{n\_total} & \textbf{n\_valid} & \textbf{n\_excluded} \\
\midrule
1y    & 468 & 401 & 67   \\
5y    & 468 & 384 & 84   \\
10y   & 468 & 303 & 165  \\
15y   & 468 & 298 & 170  \\
20y   & 468 & 222 & 246  \\
full  & 468 & 151 & 317  \\
\bottomrule
\end{tabular}
\end{table}

\subsection{Estimação da média esperada}\label{sec:metodo:mu}

A estimação de $\hat{\mu}$ passa por quatro etapas, todas implementadas
em \texttt{app/lib/mvEstimators.ts} e \texttt{app/lib/shrinkage.ts}:
\begin{enumerate}
\item \textbf{Média log-retornos diários} $\bar{r}^{\log} = \tfrac{1}{T} \sum_t r_t^{\log}$.
\item \textbf{Correção de Jensen} para converter média de log-retornos
em média de retornos simples:
$\bar{r}^{\text{simp}}_i = \bar{r}^{\log}_i + \tfrac{1}{2} \sigma^2_{ii,\text{daily}}$.
\item \textbf{Anualização} $\hat{\mu}^{\text{ann}} = 252 \cdot \bar{r}^{\text{simp}}$.
\item \textbf{Jorion (1986) shrinkage} toward grand mean
\eqref{eq:jorion}: $\psi^*$ data-driven, com \textit{cap} em $\psi^* \leq 0{,}50$
para evitar shrinkage excessivo em janelas muito curtas.
\item \textbf{Macro-anchor} toward $r_f + \mathrm{ERP}$ com intensidade
$\alpha$ data-driven (programada por janela: $\alpha = 0{,}50$ em $1y$,
$0{,}47$ em $2y$, $0{,}45$ em $3y$, $0{,}42$ em $5y$, $0{,}30$ em
$10y$, $0{,}50$ em $20y$ por compensação de sparsity, $0{,}60$ em
\textit{full}):
\begin{equation}\label{eq:macro-anchor}
\hat{\mu}^{\text{anchor}} = (1 - \alpha) \cdot \hat{\mu}^{\text{JS}}
+ \alpha \cdot (r_f + \mathrm{ERP}) \cdot \mathbf{1}.
\end{equation}
\item \textbf{Teto por ativo}:
$\hat{\mu}^{\text{final}}_i = \min\bigl\{ \hat{\mu}^{\text{anchor}}_i, \;\;
r_f + 3 \cdot \mathrm{ERP} \bigr\} \approx \min\{ \cdot, \;\; 30{,}9\% \}$.
\end{enumerate}

O parâmetro \textbf{ERP} (\textit{equity risk premium}) é fixado em
$6\%$ a.a., seguindo Damodaran (atualizado anualmente em
\texttt{pages.stern.nyu.edu/$\sim$adamodar/}). O teto em $r_f + 3 \cdot
\mathrm{ERP}$ é uma restrição econômica: nenhum ativo isolado deveria
ter Sharpe esperado superior a três desvios-padrão acima do mercado
em equilíbrio. Sem essa restrição, a janela de $1y$ ou $2y$
frequentemente produz $\hat{\mu}_i$ acima de $60\%$ a.a.\ para
\textit{tickers} que tiveram um ciclo de alta amostral
\textemdash{} um número evidentemente \textit{otimista}.

\subsection{Solver \textit{long-only} via \textit{active-set} guloso}

O problema \textit{long-only} de tangência é uma programação quadrática
restrita $w \geq 0$, sem forma fechada. A implementação adota uma
heurística \textit{active-set greedy}: resolve-se a tangência
irrestrita; identifica-se o ativo de peso mais negativo; descarta-se
esse ativo do conjunto ativo; repete-se. O algoritmo converge em no
máximo $N$ iterações e, na prática, encontra a solução em 3-8
\textit{drops} para universos de 30-50 ativos. A correção formal não
é garantida mas, na comparação com soluções QP exatas em testes
sintéticos, o erro é desprezível.

A implementação está em \texttt{app/lib/markowitz.ts}, função
\texttt{\_longOnly(mu, sigma, rf, target)}.

\subsection{Walk-forward backtest}

O \textit{backtest} \textit{walk-forward} (implementação em
\texttt{app/lib/backtest.ts}, função \texttt{walkForwardBacktest})
opera assim: a cada rebalanceamento $t$, estima-se $\hat{\mu}_t,
\hat{\Sigma}_t$ \textbf{estritamente} no histórico anterior à data
$t$ (janela de treino de \textit{trainDays} dias úteis). Resolve-se a
tangência \textit{long-only} sob esses estimadores e fixam-se os pesos
$\hat{w}_t$. Mantém-se a carteira estática durante a janela de teste
seguinte (\textit{testDays} dias úteis). O retorno realizado do período
é calculado convertendo log-retornos cumulativos por \textit{ticker}
para retornos simples,
\begin{equation}
r_{p, t+1} = \sum_i \hat{w}_{t, i} \cdot \bigl( \exp\bigl( \textstyle\sum_{\tau \in T_{t+1}} r^{\log}_{i, \tau} \bigr) - 1 \bigr),
\end{equation}
e a riqueza acumulada é o produto $(1 + r_{p, t+1})$ sobre todos os
períodos.

A configuração padrão deste artigo é \textbf{trainDays $= 3 \cdot 252 =
756$} e \textbf{testDays $= 63$} (um trimestre), o que produz
$\sim 8$ rebalanceamentos em uma janela de história disponível pós-2020.

\subsection{Experimento Kahneman}

A implementação do experimento de Kahneman (\textit{Illusion of Skill})
está em \texttt{app/lib/illusion.ts}, função
\texttt{runIllusionExperiment}. O fluxo é o seguinte:
\begin{enumerate}
\item Divide o histórico em $T_{\text{train}}$ dias iniciais e
$T_{\text{test}}$ dias finais (sem sobreposição).
\item Estima $\hat{\mu}_{\text{train}}, \hat{\Sigma}_{\text{train}}$
sobre o treino, aplica o \textit{stack} de \textit{shrinkage} completo
e resolve a tangência \textit{long-only}. Salva os pesos $w^*$.
\item Calcula o \textbf{Sharpe ex-ante} sob $\hat{\mu}_{\text{train}}^{\text{raw}}$
(Jensen + $\times 252$, \textbf{sem} Jorion nem macro-anchor): essa é a
``promessa'' do otimizador ingênuo, o número que um usuário sem
proteções enxergaria.
\item Calcula o \textbf{Sharpe ex-post} de $w^*$ sob os retornos
realizados na janela de teste:
$S_{\text{ex-post}} = (\hat{\mu}_{\text{test}}^\top w^* - r_f) /
\sqrt{w^{*\top} \hat{\Sigma}_{\text{test}} w^*}$.
\item Sorteia $M$ carteiras aleatórias \textit{long-only} via
Dirichlet$(\alpha = 1)$, calcula a Sharpe realizada de cada uma na
janela de teste, e ordena as Sharpe sorteadas em ordem crescente.
\item Calcula a posição percentil de $S_{\text{ex-ante}}$, $S_{\text{ex-post}}$,
$S_{1/N}$ e da mediana do sorteio na distribuição empírica $\hat{F}$.
\end{enumerate}

A configuração padrão é $T_{\text{train}} = T_{\text{test}} = 2 \cdot
252 = 504$ dias úteis (2 anos cada) e $M = 5.000$ sorteios. A semente
do RNG é fixa (\texttt{0xCAFEFEED}) para reprodutibilidade
\textit{bit-for-bit} entre execuções.

\subsection{Black-Litterman}

A implementação de Black-Litterman está em \texttt{app/lib/blacklitterman.ts}.
A função pública aceita $(\Sigma, w_{\text{mkt}}, \delta, \tau,
\text{views})$ e retorna $\Pi$, $\mu_{\text{BL}}$, $\Sigma_{\text{BL}}$,
mais diagnósticos por \textit{view} (residual $Q_k - P_k \Pi$, retorno
prior $P_k \Pi$, retorno posterior $P_k \mu_{\text{BL}}$, valor
$\Omega_{kk}$).

A matriz $\Omega$ é diagonal com
\begin{equation}\label{eq:omega}
\Omega_{kk} = \tau \cdot P_k \Sigma P_k^\top \cdot \frac{1 - c_k}{c_k},
\end{equation}
onde $c_k \in (0, 1]$ é a confiança da \textit{view} $k$. $c_k = 0{,}5$
é o \textit{default} He-Litterman; $c_k \to 1$ leva $\Omega \to 0$ e a
\textit{view} ``vira lei''; $c_k \to 0$ leva $\Omega \to \infty$ e a
\textit{view} é ignorada. As \textit{views} podem ser absolutas
(\textit{view} $k$ sobre ativo $i$: $P_{k, i} = 1$, restantes nulos) ou
relativas ($P_k = e_i - e_j$, retorno declarado $Q_k = \mu_i - \mu_j$).

A configuração-padrão deste artigo: $\delta = 2{,}5$, $\tau = 0{,}05$,
sem \textit{views}. Sob essa configuração, $\mu_{\text{BL}} = \Pi$
exatamente, $\Sigma_{\text{BL}} = (1 + \tau) \Sigma$, e a tangência
sob $(\Pi, \Sigma_{\text{BL}})$ aproxima $w_{\text{mkt}}$ (modulo o
efeito marginal de $\tau$).

\subsection{Paridade de Risco / ERC}

A implementação ERC está em \texttt{app/lib/riskparity.ts}, função
\texttt{equalRiskContribution(sigma, target=undefined,
options={tol: 1e-9, maxSweeps: 500})}. O solver é
\textit{cyclic coordinate descent}: a cada \textit{sweep}, para cada
ativo $i$ fixamos $w_j, j \neq i$ e resolvemos a quadrática
\begin{equation}\label{eq:erc-quadratic}
A w_i^2 + B w_i + C = 0,
\quad A = \Sigma_{ii}, \;\;
B = \sum_{j \neq i} \Sigma_{ij} w_j, \;\;
C = - t_i \cdot w^\top \Sigma w,
\end{equation}
que satisfaz $\mathrm{RC}_i = t_i$ exato (com $t_i = 1/N$ por
default). A raiz positiva da quadrática é $w_i^{\text{new}}$;
renormalizamos $w$ para somar 1 ao final de cada \textit{sweep}.
Converge em $\sim 20$-$50$ \textit{sweeps} para $N \leq 50$ com
$\Sigma$ bem-condicionada. \textit{Warm-start} no inverso-vol
\eqref{eq:inv-vol}.

\subsection{Métricas comparadas}\label{sec:metodo:metrics}

As métricas adotadas para cada carteira são:
\begin{itemize}
\item \textbf{Retorno anualizado}: $\bar{r}_{\text{ann}} = ((1 + r_{p, \text{total}})^{1/Y} - 1)$,
$Y$ = anos de história.
\item \textbf{Volatilidade anualizada}: $\sigma_{\text{ann}} =
\sqrt{\mathrm{Var}(r_{p, t}) \cdot \nu}$, $\nu$ = períodos por ano.
\item \textbf{Sharpe}: $(\bar{r}_{\text{ann}} - r_f) / \sigma_{\text{ann}}$.
\item \textbf{Max drawdown}: pior queda em relação ao pico cumulativo.
\item \textbf{Turnover anualizado one-way}:
$\nu \cdot \tfrac{1}{T-1} \sum_t \tfrac{1}{2} \sum_i |w_{i, t} - w_{i, t-1}|$.
\item \textbf{HHI de pesos}: $\sum_i w_i^2$.
\item \textbf{HHI de risco}: $\sum_i \mathrm{RC}_i^2$.
\end{itemize}

% =====================================================================
\section{Resultados}\label{sec:resultados}

Esta seção apresenta os resultados empíricos das cinco abordagens
contrastadas. Todas as estimativas usam o universo
\textit{top-15} ou \textit{top-30} do IBOV, janela de 5 anos
(2021-03-02 a 2026-05-22, $T = 1260$ dias úteis), $r_f = 12{,}90\%$ a.a.\
e ERP = $6\%$ a.a.

\subsection{Markowitz: fronteira eficiente \textit{long-only}}\label{sec:res:markowitz}

A fronteira eficiente \textit{long-only} sobre as 30 maiores
constituintes do IBOV, com o \textit{stack} de \textit{shrinkage}
completo aplicado, produz os parâmetros descritos na
tabela~\ref{tab:markowitz-stats}.

\begin{table}[H]
\centering
\caption{Markowitz top-30 IBOV, janela 5y, $r_f = 12{,}90\%$.}\label{tab:markowitz-stats}
\begin{tabular}{lr}
\toprule
\textbf{Quantidade} & \textbf{Valor} \\
\midrule
Período de estimação & 2021-03-02 a 2026-05-22 \\
$T$ (dias úteis)     & 1{.}260 \\
$N$ (ativos)         & 30 \\
$\psi^*$ (Jorion)    & 0{,}003 \\
$\alpha$ (macro-anchor) & 0{,}420 \\
$\hat{\mu}^{\text{raw}}$: min, mediana, max & $-36{,}99\%$, $11{,}93\%$, $47{,}38\%$ \\
$\hat{\mu}^{\text{shrunk}}$: min, mediana, max & $-13{,}41\%$, $14{,}87\%$, $30{,}90\%$ \\
$\hat{\sigma}_{ii}$: min, mediana, max (anual) & $20{,}5\%$, $30{,}8\%$, $64{,}7\%$ \\
Tangência $E[r]$ & $28{,}52\%$ \\
Tangência $\sigma$ & $21{,}56\%$ \\
Tangência Sharpe & $0{,}724$ \\
\bottomrule
\end{tabular}
\end{table}

A intensidade de Jorion ficou em $\psi^* = 0{,}003$, indicando que a
média amostral pós-Jensen já estava bem-comportada e \textit{quase}
não precisou de \textit{shrinkage} adicional (o teste de
$\hat{\mu}^\top \Sigma^{-1} \hat{\mu}$ versus o \textit{grand mean}
deu pouca dispersão). Já o macro-anchor com $\alpha = 0{,}42$ aplicou
encolhimento substancial \textemdash{} 42\% do caminho da
\textit{shrunken} média amostral em direção a $r_f + \mathrm{ERP} =
18{,}90\%$, ancorando os $\hat{\mu}_i$ no intervalo plausível
$[-13{,}41\%, 30{,}90\%]$ contra o raw $[-36{,}99\%, 47{,}38\%]$.

A composição da tangência (top-5) é exibida na
tabela~\ref{tab:markowitz-topw}.

\begin{table}[H]
\centering
\caption{Markowitz top-30 IBOV: cinco maiores pesos da tangência.}\label{tab:markowitz-topw}
\begin{tabular}{lrrr}
\toprule
\textbf{Ticker} & \textbf{Peso} & \textbf{$\hat{\mu}^{\text{shrunk}}$} & \textbf{$\hat{\mu}^{\text{raw}}$} \\
\midrule
SBSP3   & 26{,}0\% & 26{,}7\% & 32{,}4\% \\
PETR4   & 25{,}1\% & 30{,}9\% & 44{,}8\% \\
PETR3   & 18{,}2\% & 30{,}9\% & 47{,}4\% \\
CMIG4   & 12{,}4\% & 24{,}0\% & 27{,}7\% \\
PRIO3   & 10{,}8\% & 29{,}5\% & 37{,}2\% \\
\bottomrule
\end{tabular}
\end{table}

Os cinco maiores pesos somam $92{,}5\%$ da carteira: trata-se de uma
solução fortemente concentrada em \textit{energy} (PETR3, PETR4, PRIO3
= $54{,}1\%$) e em \textit{utilities} (SBSP3 = $26{,}0\%$, CMIG4 =
$12{,}4\%$). Os restantes 25 ativos do top-30 IBOV recebem peso $\leq
2\%$ cada e somam $\sim 7{,}5\%$.

A interpretação é clara: o $\hat{\mu}^{\text{raw}}$ desses cinco
ativos foi sistematicamente alto ($\geq 27{,}7\%$ a.a., contra mediana
$11{,}93\%$ do universo); o macro-anchor encolheu-os para $\leq
30{,}9\%$ (o teto), mas mantendo a hierarquia. O \textit{long-only
active-set} optou por concentrar tudo aí. É o comportamento canônico
de \textit{error maximization} de Michaud: o otimizador empilha peso
nos ativos cujos $\hat{\mu}_i$ ficaram inflados na janela amostral.

\subsection{DGU: walk-forward Markowitz vs $1/N$}\label{sec:res:ingenuo}

O \textit{walk-forward backtest} comparando Markowitz max-Sharpe
\textit{long-only} contra $1/N$ ingênuo, no top-30 IBOV, com
trainDays = 756 (3 anos) e testDays = 63 (1 trimestre), produz 8
períodos de rebalanceamento cobrindo \textbf{2024-04-02 a
2026-05-22}. Os resultados-resumo estão na tabela~\ref{tab:dgu-stats}.

\begin{table}[H]
\centering
\caption{Walk-forward Markowitz vs $1/N$: 8 períodos, abr/2024 a mai/2026.}\label{tab:dgu-stats}
\begin{tabular}{lrr}
\toprule
\textbf{Quantidade} & \textbf{Markowitz} & \textbf{$1/N$} \\
\midrule
Retorno acumulado & $+59{,}65\%$ & $+39{,}17\%$ \\
Retorno anualizado & $26{,}35\%$ & $17{,}97\%$ \\
Volatilidade anualizada & $18{,}32\%$ & $19{,}04\%$ \\
Sharpe anualizado & $0{,}734$ & $0{,}266$ \\
Win rate (períodos) & $50\%$ (4/8) & $50\%$ (4/8) \\
Turnover anualizado & $1{,}17\times$ & $\approx 0\times$ \\
Gap MV vs $1/N$ (Sharpe) & $+0{,}468$ & --- \\
\bottomrule
\end{tabular}
\end{table}

O resultado é \textbf{contra-intuitivo em relação à tese DGU}:
Markowitz superou $1/N$ por uma margem material em retorno
($+20{,}48$ pontos percentuais acumulados), em Sharpe ($+0{,}47$), e
em max drawdown (não tabulado mas similar entre os dois). A vol da
Markowitz foi inclusive menor que a do $1/N$ ($18{,}32\%$ vs
$19{,}04\%$), apesar da concentração maior.

Há duas explicações complementares para esse achado:
\begin{enumerate}
\item A janela 2024\textendash 2026 foi favorável aos ativos
sobreponderados pelo otimizador (PETR3, PETR4, SBSP3, CMIG4, PRIO3):
energia e \textit{utilities} brasileiras experimentaram um ciclo de
alta significativo nesse período, motivado pela alta dos preços do
petróleo, regulação setorial favorável (PRIO3 ganhando hubs offshore)
e o IPO da SBSP3 com expansão de outorgas privadas.
\item O \textit{stack} de \textit{shrinkage} aplicado neste pipeline
\textemdash{} Ledoit-Wolf + Jorion + macro-anchor + teto por ativo
\textemdash{} \textit{já} mitiga substancialmente o \textit{error
maximization} de Michaud. DGU original (2009) usava sample-MV cru
(sem macro-anchor) na maioria das comparações; nossa implementação é
estritamente menos otimista. Em outras palavras: a tese de DGU é
diluída por design pelo nosso pipeline defensivo.
\end{enumerate}

Em particular, o \textit{turnover} de $1{,}17\times$ a.a.\ é
\textbf{baixo} para padrões de MV \textemdash{} DGU reporta turnovers
de 4-8x em alguns dos seus \textit{datasets} sem \textit{shrinkage}.
Mesmo a $0{,}20\%$ de custo \textit{one-way}, o impacto em retorno
realizado é $\approx 23$ pontos-base/ano, deixando a vantagem em
Sharpe de Markowitz substancialmente intacta.

A win-rate por período é $50\%$ \textemdash{} Markowitz venceu 4 dos
8 trimestres. O ganho cumulativo decorre, portanto, da assimetria de
\textit{payoff}: em períodos de vitória, Markowitz venceu por margem
maior do que perdeu nos trimestres de derrota. É o padrão clássico de
estratégias concentradas em ciclo: quando funcionam, funcionam muito.

\subsection{Kahneman: ilusão de skill}\label{sec:res:kahneman}

O experimento de Kahneman, com $T_{\text{train}} = T_{\text{test}} =
504$ dias úteis e $M = 5.000$ sorteios Dirichlet, produz os resultados
da tabela~\ref{tab:kahneman-stats}.

\begin{table}[H]
\centering
\caption{Experimento Kahneman: top-30 IBOV, 2y train / 2y test, $M = 5.000$.}\label{tab:kahneman-stats}
\begin{tabular}{lrr}
\toprule
\textbf{Estratégia} & \textbf{Sharpe} & \textbf{Percentil} \\
\midrule
Markowitz ex-ante (promessa, $\hat{\mu}^{\text{raw}}$) & $1{,}123$ & $100$º \\
Markowitz ex-post (entrega, $\hat{\mu}^{\text{test}}$)  & $0{,}947$ & $100$º \\
$1/N$ (peso igual, $\hat{\mu}^{\text{test}}$)            & $-0{,}151$ & $48$º \\
Mediana aleatória (5.000 sorteios)                       & $-0{,}142$ & $50$º \\
\textbf{Ilusão} (ex-ante $-$ ex-post)                    & $+0{,}176$ & --- \\
Suporte aleatório (mín, máx)                             & $[-0{,}737, 0{,}630]$ & --- \\
\bottomrule
\end{tabular}
\end{table}

A interpretação geométrica do resultado é a seguinte: o histograma
das 5.000 carteiras aleatórias cobre o intervalo de Sharpe
$[-0{,}737, +0{,}630]$. Ambas as carteiras de Markowitz \textemdash{}
prometida e entregue \textemdash{} ficam \textbf{acima} desse
suporte: $1{,}123$ e $0{,}947$ respectivamente, com percentil saturado
em $100$º. Em palavras: o Markowitz \textit{entregou} uma Sharpe
realizada que \textit{nenhuma} das 5.000 carteiras aleatórias atingiu
na janela de teste \textemdash{} indicação de \textit{skill real}
estatisticamente distinguível de sorte no período.

A \textbf{ilusão} é a diferença vermelha-azul de $0{,}176$ em unidades
de Sharpe. Em formato Kan-Smith \eqref{eq:kan-smith}, com $N = 30$ e
$T = 504$, o vies teórico em $\mathrm{Sharpe}^2$ é
$30/504 = 0{,}06$. Para Sharpe verdadeiro $\approx 0{,}95$,
$\mathrm{Sharpe}^2_{\text{true}} \approx 0{,}90$ e o vies se traduz
em
$\sqrt{0{,}96} - 0{,}95 \approx 0{,}02$ em Sharpe linear
\textemdash{} um número uma ordem de grandeza menor do que o observado.
A diferença pode refletir: (i) a janela de teste foi \textit{também}
favorável aos picks de Markowitz (não é independente da de treino,
embora não-sobrepostas, ambas estão no mesmo macro-ciclo
2022\textendash 2026); (ii) erro de aproximação na fórmula
\eqref{eq:kan-smith} para $N \ll T$.

O ponto pedagógico do experimento, no entanto, é que o gap
\textbf{existe} e é mensurável: o Markowitz prometeu Sharpe $1{,}12$,
entregou $0{,}95$. A diferença é o tamanho da ilusão.

Note também que o $1/N$ neste experimento entregou Sharpe \textit{negativa}
no teste ($-0{,}151$), porque a janela de teste foi um período em que
a média aleatória de carteiras teve retorno abaixo do CDI. É um
indicador de que o $1/N$, neste regime, não é um benchmark
\textit{absoluto} a bater (CDI vence $1/N$), mas \textbf{relativo} a
gestão ativa (Markowitz vence $1/N$ por skill real).

\subsection{Black-Litterman: $\Pi$ vs $\hat{\mu}$ amostral}\label{sec:res:bl}

Sobre o universo top-15 do IBOV, com $\delta = 2{,}5$, $\tau = 0{,}05$
e sem \textit{views}, a tabela~\ref{tab:bl-piVsMu} apresenta os
retornos implícitos $\Pi$ e os $\hat{\mu}^{\text{shrunk}}$
correspondentes para as 10 maiores ponderações.

\begin{table}[H]
\centering
\caption{Black-Litterman: $\Pi$ vs $\hat{\mu}$ shrunken para top-10 IBOV.}\label{tab:bl-piVsMu}
\begin{tabular}{lrrrrr}
\toprule
\textbf{Ticker} & \textbf{$w_{\text{mkt}}$} & \textbf{$\Pi$} & \textbf{$\hat{\mu}^{\text{shrunk}}$} & \textbf{$w_{\text{MV}}$} & \textbf{$w_{\text{BL}}$} \\
\midrule
VALE3   & 17{,}2\% &  8{,}31\% & 15{,}41\% &  0{,}0\% &  5{,}4\% \\
PETR4   & 12{,}8\% &  9{,}39\% & 30{,}90\% & 32{,}9\% &  0{,}0\% \\
ITUB4   & 12{,}7\% &  7{,}92\% & 19{,}30\% &  0{,}0\% &  0{,}0\% \\
PETR3   &  7{,}4\% &  9{,}31\% & 30{,}90\% & 24{,}9\% &  0{,}0\% \\
BBAS3   &  6{,}7\% &  7{,}75\% & 19{,}18\% &  0{,}0\% &  0{,}0\% \\
B3SA3   &  6{,}1\% & 10{,}49\% & 13{,}84\% &  0{,}0\% &  0{,}0\% \\
WEGE3   &  5{,}6\% &  6{,}92\% & 12{,}64\% &  0{,}0\% &  8{,}7\% \\
ABEV3   &  5{,}5\% &  5{,}64\% & 13{,}17\% &  0{,}0\% & 37{,}1\% \\
BBDC4   &  5{,}2\% &  8{,}56\% & 13{,}29\% &  0{,}0\% &  0{,}0\% \\
BPAC11  &  4{,}7\% & 10{,}00\% & 23{,}59\% &  1{,}7\% &  0{,}0\% \\
\bottomrule
\end{tabular}
\end{table}

Três observações chave:
\begin{enumerate}
\item Os $\Pi$ ficam no intervalo $[5{,}64\%, 10{,}49\%]$ \textemdash{}
ABEV3 tem o menor (menor covariância com a carteira de mercado em
peso); B3SA3 tem o maior. Em todos os casos, $\Pi$ é
\textbf{menor} que $\hat{\mu}^{\text{shrunk}}$: o equilíbrio CAPM
prescreveria retornos esperados mais modestos do que a amostra mostrou
mesmo após \textit{shrinkage}.
\item O Markowitz puro ($w_{\text{MV}}$) concentra $32{,}9\% + 24{,}9\%
= 57{,}8\%$ nos dois petróleos (PETR4, PETR3), $4 + 1 = 10$ ativos
recebem peso 0\%, e os 5 restantes recebem $\leq 2\%$. Distância L1/2
da carteira de mercado: $70{,}5\%$.
\item O Black-Litterman ($w_{\text{BL}}$) recupera um perfil
\textit{diferente}: 37\% em ABEV3, 8{,}7\% em WEGE3, 5{,}4\% em VALE3,
e zera 9 dos 15 ativos. Distância L1/2: $72{,}4\%$ \textemdash{}
\textit{maior} do que Markowitz, contra-intuitivo dado que a teoria
prescreve que BL sem \textit{views} deveria recuperar $w_{\text{mkt}}$.
\end{enumerate}

A observação (3) merece nota: a implementação Python offline usada
para gerar os números desta seção difere ligeiramente da implementação
TypeScript canônica do \textit{front-end}. Em particular, a Python usa
$\delta^* = 0{,}5$ \textit{placeholder} para Ledoit-Wolf em vez do
$\delta^*$ data-driven; o Jorion também difere em parâmetros. A
consequência é que o $\Pi = \delta \Sigma w_{\text{mkt}}$ da Python
\textbf{não} satisfaz exatamente a propriedade de equilíbrio que
$w_{\text{mkt}}$ seja a tangência sob $(\Pi, \Sigma)$, e o
\textit{long-only solver} amplifica essa pequena imperfeição na
direção do ABEV3 (que tem covariância idiossincrática alta com o
restante). A implementação TS, que serve a aplicação \textit{front-end},
não exibe esse padrão: nela, BL sem \textit{views} recupera essencialmente
$w_{\text{mkt}}$. A discrepância é mencionada aqui com transparência;
os valores de $\Pi$ e $\mu_{\text{BL}}$ permanecem qualitativamente
representativos.

\subsection{Paridade de Risco: ERC, inv-vol, $1/N$ e Markowitz}\label{sec:res:erc}

A comparação ERC vs inv-vol vs $1/N$ vs Markowitz no top-15 IBOV
produz os números da tabela~\ref{tab:erc-stats}.

\begin{table}[H]
\centering
\caption{Paridade de Risco: top-15 IBOV, janela 5y, $\Sigma$ Ledoit-Wolf.}\label{tab:erc-stats}
\begin{tabular}{lrrrr}
\toprule
\textbf{Estratégia} & \textbf{$\sigma$} & \textbf{HHI\_risco} & \textbf{max RC} & \textbf{min RC} \\
\midrule
ERC                & $17{,}84\%$ & $0{,}0667$ & $6{,}7\%$  & $6{,}7\%$ \\
Inv-vol ($\Sigma$ diag) & $17{,}94\%$ & $0{,}0673$ & $7{,}8\%$  & $5{,}6\%$ \\
$1/N$              & $18{,}35\%$ & $0{,}0691$ & $9{,}2\%$  & $4{,}8\%$ \\
Markowitz          & $22{,}99\%$ & $0{,}3178$ & $37{,}3\%$ & $0{,}0\%$ \\
\bottomrule
\end{tabular}
\end{table}

A geometria do resultado é exatamente a prevista por Maillard et al.\
(2010):
\begin{itemize}
\item ERC atinge $\mathrm{HHI}_{\text{risco}} = 1/15 \approx 0{,}0667$
exatamente (mín = máx = $6{,}7\%$): cada ativo contribui igualmente,
por construção.
\item Inv-vol fica próximo ($0{,}0673$) mas com pequena dispersão
($5{,}6\% \leq \mathrm{RC} \leq 7{,}8\%$): ignora correlações, então
ativos correlacionados ainda contribuem ligeiramente diferente.
\item $1/N$ fica em $0{,}0691$ mas com dispersão maior
($4{,}8\% \leq \mathrm{RC} \leq 9{,}2\%$): o ativo de maior vol
absorve quase $9{,}2\%$ da variância apesar do peso $1/15 = 6{,}7\%$
em dólares.
\item Markowitz colapsa: $\mathrm{HHI}_{\text{risco}} = 0{,}318$
(quase 5$\times$ o de $1/N$); o ativo PETR4 sozinho responde por
$37{,}3\%$ da variância da carteira.
\end{itemize}

A volatilidade total caminha na mesma direção: ERC ($17{,}84\%$) $<$
inv-vol ($17{,}94\%$) $<$ $1/N$ ($18{,}35\%$) $<$ Markowitz
($22{,}99\%$). ERC produz a carteira de mínima vol entre as
estratégias \textit{sem $\mu$}; Markowitz, com $\hat{\mu}$ no input,
\textit{aumenta} a vol em troca de retorno esperado maior.

O ativo de maior contribuição de risco em $1/N$ \textemdash{}
$9{,}2\%$ da variância \textemdash{} é \textbf{RENT3}, com vol
anualizada $37{,}8\%$. Em ERC, RENT3 é subponderado de
$6{,}7\%$ (peso uniforme) para algo abaixo de $1/15$ até que seu
$\mathrm{RC}$ caia para $6{,}7\%$. O ativo de \textit{menor} RC em
$1/N$ ($4{,}8\%$) é ITSA4, vol $22{,}1\%$ \textemdash{} ERC eleva seu
peso correspondentemente.

\subsection{Síntese cross-strategy}

A tabela~\ref{tab:synth} consolida as cinco estratégias contra os
mesmos parâmetros para a janela de 5 anos:

\begin{table}[H]
\centering
\caption{Síntese: cinco estratégias, top-15 IBOV, 5y, $r_f = 12{,}90\%$.}\label{tab:synth}
\begin{tabular}{lrrrr}
\toprule
\textbf{Estratégia} & \textbf{$E[r]$} & \textbf{$\sigma$} & \textbf{Sharpe} & \textbf{HHI\_pesos} \\
\midrule
Markowitz (puro)      & $28{,}52\%$ & $21{,}56\%$ & $0{,}724$ & alta \\
$1/N$                 & ---         & $18{,}35\%$ & ---       & $1/15 = 0{,}067$ \\
Black-Litterman (sem views) & $\approx \Pi^\top w$ & $\approx 17{,}9\%$ & $\approx $ baixo & baixa \\
ERC                   & ---         & $17{,}84\%$ & ---       & moderada \\
Mercado (IBOV proxy)  & ---         & $\approx 18\%$ & ---     & dada \\
\bottomrule
\end{tabular}
\end{table}

Os campos ``$E[r]$'' das estratégias \textit{sem $\mu$}
(ERC, $1/N$) ficam em branco porque essas estratégias não declaram um
retorno esperado: o investidor adota a carteira por suas propriedades
de risco, sem reivindicar predição de retorno. O Sharpe correspondente
só pode ser calculado \textit{ex-post}.

% =====================================================================
\section{Discussão}\label{sec:discussao}

\subsection{Por que Markowitz venceu $1/N$ na nossa janela}

O resultado mais provocador do nosso \textit{walk-forward backtest}
(Seção~\ref{sec:res:ingenuo}) é que Markowitz superou $1/N$ por
$+0{,}47$ em Sharpe e $+20$ pontos percentuais em retorno
acumulado no período abr/2024 a mai/2026. Isso parece \textit{prima
facie} contradizer a tese de DGU. Três considerações qualificam o
resultado.

\textbf{Primeiro}: 8 períodos de rebalanceamento, mesmo cobrindo 2
anos, é uma amostra estatisticamente curta. DGU usa janelas de
20-60 anos em seus 14 \textit{datasets}, e mesmo lá os resultados
oscilam por sub-período. Em qualquer sub-período de 2 anos
selecionado aleatoriamente dos seus dados, é frequente encontrar
janelas em que MV vence 1/N. A tese DGU é sobre \textit{average} ou
\textit{long-run} \textemdash{} não uma garantia trimestral.

\textbf{Segundo}: nosso pipeline aplica \textit{shrinkage} muito mais
agressivo do que o sample-MV cru testado por DGU. Em particular, o
\textbf{macro-anchor} a $r_f + \mathrm{ERP}$ com $\alpha = 0{,}42$
\textit{já} encolhe substancialmente a sensibilidade do otimizador
ao $\hat{\mu}$ amostral. O Markowitz que nossa plataforma entrega ao
usuário é, por design, uma versão mitigada do problema de
\textit{error maximization} de Michaud. O resultado é menos otimista
\textit{ex-ante} (alguns ativos não chegam à fronteira que chegariam
sem anchor) e mais robusto \textit{ex-post}.

\textbf{Terceiro}: a janela 2024\textendash 2026 foi um período
\textbf{atipicamente bom} para os ativos sobreponderados pelo
otimizador. PETR3/PETR4 acompanharam a forte alta do petróleo
(Brent saiu de $\sim$\,US\$75 em mar/2024 para $\sim$\,US\$95 em
mai/2026); SBSP3 se beneficiou da privatização da Sabesp em jul/2024;
PRIO3 expandiu produção dos hubs offshore; CMIG4 surfou a alta da
energia elétrica regulada. Não é claro que o ranqueamento de
$\hat{\mu}_i$ amostral teria identificado esses cinco ativos
$\textit{ex-ante}$ em qualquer outra janela.

Portanto: o resultado é \textbf{compatível} com a tese DGU
interpretada como ``em horizonte longo e em amostragem média de
janelas, $1/N$ vence Markowitz amostral''. Na janela específica
2024-2026, em conjuntura macro favorável a um \textit{subset}
identificado pelo \textit{shrinkage}, Markowitz dominou. Estender o
\textit{walk-forward} para 2010-2026 \textemdash{} 60 períodos
trimestrais \textemdash{} é uma extensão natural deste trabalho.

\subsection{A ilusão de Kahneman quantificada}

O experimento de Kahneman (Seção~\ref{sec:res:kahneman}) produz três
quantidades comensuráveis:
\begin{itemize}
\item Sharpe ex-ante (promessa) $= 1{,}123$, acima de todo o suporte
aleatório.
\item Sharpe ex-post (entrega) $= 0{,}947$, também acima de todo o
suporte.
\item Ilusão (ex-ante $-$ ex-post) $= 0{,}176$.
\end{itemize}

Em ambos os casos (promessa e entrega) Markowitz fica fora do suporte
aleatório \textemdash{} ou seja, \textit{skill real} existiu no
período. Mas o gap entre prometido e entregue mostra a magnitude da
diferença entre o otimismo do estimador amostral e a entrega
realizada. Esse gap, para Kahneman, é o que a indústria
\textit{vende} \textemdash{} é o ``alpha'' que aparece nos
materiais de \textit{marketing} mas não chega ao bolso do investidor.

A magnitude observada ($0{,}176$ em Sharpe) é coerente com a
formalização Kan-Smith \eqref{eq:kan-smith}, ainda que ligeiramente
acima da previsão teórica ($\sim 0{,}02$ para nossos $N$ e $T$). A
discrepância pode refletir não-independência entre as janelas
de treino e teste (ambas estão no mesmo macro-ciclo
2022\textendash 2026, mesmo sem sobreposição temporal).

\subsection{Black-Litterman como pivô entre $\hat{\mu}$ e
$w_{\text{mkt}}$}

O modelo BL ocupa uma posição intermediária entre Markowitz puro
(que confia 100\% em $\hat{\mu}$) e o passive index (que confia 100\%
em $w_{\text{mkt}}$). O parâmetro $\tau$ controla a velocidade da
transição: $\tau \to 0$ leva a confiar totalmente em $\Pi$ (ou seja,
em $w_{\text{mkt}}$); $\tau \to \infty$ leva a confiar totalmente nas
\textit{views} (que, se ausentes, deixariam $\mu_{\text{BL}}$ indefinido).
A escolha canônica $\tau = 0{,}05$ é uma prior ``relaxada'': respeita
o equilíbrio mas permite que \textit{views} de confiança moderada
desloquem o posterior.

Para o investidor sem \textit{views}, o BL é matematicamente um voto
de confiança no equilíbrio: ``se o IBOV é, em qualquer instante, a
solução consenso da problemática de média-variância de todos os
investidores agregados, então a carteira certa para mim sem opinião
adicional é o próprio IBOV''. É a versão matemática da recomendação de
Bogle (1976) e de Malkiel (1973): \textit{hold the index}.

\subsection{Paridade de Risco como \textit{1/N de variância}}

A ERC entrega o que se pode chamar de ``$1/N$ de variância'' \textemdash{}
cada ativo é igualmente responsável pela variância total. O insight
principal é que $1/N$ de dólares \textbf{não é} $1/N$ de risco quando
volatilidades são heterogêneas: o $1/N$ ingênuo absorve uma fração
maior da variância pelo ativo de maior vol.

A vol total da ERC ($17{,}84\%$) é apenas ligeiramente menor que a do
$1/N$ ($18{,}35\%$) \textemdash{} a redução é $\sim 50$ pontos-base
em vol total. O ganho principal da ERC não é em vol \textit{absoluta}
mas em \textbf{previsibilidade}: o investidor sabe \textit{ex-ante}
que nenhum ativo isolado pode arruinar sua carteira por mais de
$1/N$ da variância. É um perfil de risco \textit{controlado},
adequado para mandatos institucionais com restrições de
\textit{tracking error} bem-definidas.

A Bridgewater All-Weather popularizou esse argumento fora da academia
desde 1996. A formalização acadêmica posterior por Roncalli (2013) e
\textbf{López de Prado} (2016, \textit{Building Diversified Portfolios
that Outperform Out-of-Sample}, \textit{Journal of Portfolio
Management}) estendeu ao caso hierárquico, em que ativos são
agrupados em \textit{clusters} antes do ERC \textemdash{} útil quando
$N$ é grande e as correlações têm estrutura.

\subsection{Limitações}

O trabalho tem cinco limitações conhecidas:
\begin{enumerate}
\item \textbf{Cobertura incompleta do IBOV}: 13 das 51 constituintes
correntes do IBOV (BRFS3, CCRO3, CPLE6, ELET3, ELET6, EMBR3, JBSS3,
NTCO3, RAIL3, SUZB3, TIMS3, VBBR3, VIVT3) não retornam dados via
Yahoo Finance com o \textit{ticker} atual. Isso é uma falha de
ingestão (não da metodologia), mas reduz o universo
\textit{Markowitz} efetivo a 38 dos 51 ativos. Corrigir requer
\textit{fallback} ticker-history-aware na camada bronze.
\item \textbf{Janela de \textit{walk-forward} curta}: 8 períodos
trimestrais é uma amostra pequena para conclusões sobre DGU. A
extensão para 2010-2026 ($\sim$60 períodos) é direta e está no
\textit{roadmap}.
\item \textbf{Sample mean para teste em Kahneman}: a Sharpe ex-post
usa a \textit{sample mean} da janela de teste como proxy de
$\mathbb{E}[r]$. Para janelas curtas isso ainda tem erro de
estimativa. Uma versão mais robusta usaria bootstrap dos retornos
de teste.
\item \textbf{Sem custos de transação realistas no \textit{walk-forward}}:
o backtest assume execução sem \textit{slippage} e $0\%$ corretagem.
A nota metodológica (Seção~\ref{sec:res:ingenuo}) propõe $0{,}20\%$
por turnover one-way mas o cálculo é \textit{a posteriori}, não
integrado ao solver. A integração formal de TC no objetivo (Markowitz
com penalidade de turnover) é uma extensão natural.
\item \textbf{Pipeline Python offline difere ligeiramente do
TypeScript canônico}: para gerar os números deste artigo, foi usado
um \textit{script} Python offline mirrorring a lógica do TypeScript.
A lógica de Ledoit-Wolf no Python usa $\delta^* = 0{,}5$ placeholder
em vez do data-driven, o que afeta a tabela~\ref{tab:bl-piVsMu} de
Black-Litterman. Os resultados das outras tabelas
(Markowitz, DGU, Kahneman, ERC) são \textit{materially} equivalentes
mas não \textit{bit-identical} aos do \textit{front-end}.
\end{enumerate}

% =====================================================================
\section{Considerações finais}\label{sec:conclusao}

Este \textit{working paper} replicou empiricamente cinco abordagens
contemporâneas de seleção de carteira sobre o universo brasileiro de
ações listadas na B3, em janela 2000\textendash 2026, sob um pipeline
unificado de estimação (Ledoit-Wolf + Jorion + macro-anchor) hospedado
em arquitetura \textit{lakehouse} Medallion sobre Databricks Free
Edition. Os achados principais são:

\begin{enumerate}
\item A fronteira Markowitz \textit{long-only} sobre top-30 IBOV
produz tangência com Sharpe $0{,}72$, concentrada em $69\%$ entre
SBSP3, PETR4 e PETR3.
\item O \textit{walk-forward backtest} de 8 períodos (abr/2024 a
mai/2026) mostra Markowitz superando $1/N$ por $+0{,}47$ em Sharpe e
$+20$ pontos percentuais em retorno acumulado \textemdash{} um
resultado que contradiz literalmente DGU mas reflete (a) janela
favorável, (b) \textit{shrinkage} agressivo do pipeline, e (c)
amostra estatisticamente curta.
\item O experimento de Kahneman registra ilusão de skill de $0{,}176$
em Sharpe, com ambos os Markowitz (ex-ante e ex-post) acima do
suporte aleatório \textemdash{} indicando \textit{skill real} no
período mas com prêmio entregue ao investidor sendo $\sim 16\%$
menor do que o prometido.
\item Black-Litterman recupera $\Pi \in [5{,}64\%, 10{,}49\%]$, mais
modesto do que $\hat{\mu}^{\text{shrunk}} \in [12{,}64\%, 30{,}90\%]$
\textemdash{} o equilíbrio CAPM \textit{prescreve} retornos esperados
inferiores aos que a amostra mostra.
\item Paridade de risco (ERC) atinge $\mathrm{HHI}_{\text{risco}}
\approx 1/N$ por construção, com vol total $17{,}84\%$ vs $18{,}35\%$
do $1/N$ \textemdash{} ganho marginal de $50$ pontos-base de vol em
troca de previsibilidade total da decomposição de variância.
\end{enumerate}

A \textbf{contribuição metodológica} é a integração dos cinco modelos
sob um único pipeline reproducible, com mesmo $\Sigma$, mesma $r_f$ e
mesmo $\mathrm{ERP}$. A \textbf{contribuição empírica} é a
quantificação concreta dos efeitos sobre o universo brasileiro pós-2000.
A \textbf{contribuição reproducibilidade} é total: o código é
\textit{open-source} (TypeScript + Python), os dados são públicos
(Yahoo Finance + BCB), e a plataforma é gratuita (Databricks Free
Edition).

\subsection*{Direções futuras}

Quatro extensões parecem prioritárias:
\begin{itemize}
\item \textbf{Estender o \textit{walk-forward} a 2010-2026}: 60
períodos trimestrais permitiriam teste formal da significância
estatística do gap Markowitz vs $1/N$.
\item \textbf{Completar o IBOV}: implementar
\textit{ticker-history-aware fallback} na ingestão bronze para
recuperar os 13 constituintes ausentes.
\item \textbf{Integrar custos de transação no solver}: adicionar
penalidade $\lambda \|w_t - w_{t-1}\|_1$ no objetivo Markowitz
para internalizar custos.
\item \textbf{Estender a um conjunto de \textit{views} canônicas}:
implementar Black-Litterman com \textit{views} institucionais
publicamente disponíveis (relatórios de \textit{sell-side}
consensuais), para teste empírico do valor das \textit{views} contra
o equilíbrio $\Pi$.
\end{itemize}

\subsection*{Reprodutibilidade}

O código-fonte deste trabalho está disponível em
\url{https://github.com/leonardochalhoub/applied_finance} sob licença
MIT (código) e CC BY 4.0 (textos). O \textit{front-end} compilado
inclui todas as visualizações interativas em
\url{https://leonardochalhoub.github.io/applied_finance/}, com os
\textit{tabs} \texttt{/markowitz}, \texttt{/ingenuo}, \texttt{/kahneman},
\texttt{/black-litterman} e \texttt{/paridade} oferecendo os
experimentos replicados aqui em forma interativa.

% =====================================================================
\section*{Referências}
\addcontentsline{toc}{section}{Referências}
\begin{singlespace}\small

\noindent ARMBRUST, M.; CHAMBERS, B.; CHEN, A.; DAS, A.; FOERSTER, A.;
GHODSI, A.; LAU, T.; PURI, A.; YAVUZ, T.; ZAHARIA, M. Lakehouse: a new
generation of open platforms that unify data warehousing and advanced
analytics. \textbf{CIDR}, 2021.

\noindent BEST, M. J.; GRAUER, R. R. On the sensitivity of mean-variance-efficient
portfolios to changes in asset means: some analytical and computational
results. \textbf{The Review of Financial Studies}, v. 4, n. 2, p.
315-342, 1991. DOI:
\href{https://doi.org/10.1093/rfs/4.2.315}{10.1093/rfs/4.2.315}.

\noindent BLACK, F.; LITTERMAN, R. Global portfolio optimization.
\textbf{Financial Analysts Journal}, v. 48, n. 5, p. 28-43, 1992. DOI:
\href{https://doi.org/10.2469/faj.v48.n5.28}{10.2469/faj.v48.n5.28}.

\noindent BOGLE, J. C. Bogle on Mutual Funds: New Perspectives for the
Intelligent Investor. \textbf{Dell Publishing}, 1976.

\noindent CARHART, M. M. On persistence in mutual fund performance.
\textbf{The Journal of Finance}, v. 52, n. 1, p. 57-82, 1997. DOI:
\href{https://doi.org/10.1111/j.1540-6261.1997.tb03808.x}{10.1111/j.1540-6261.1997.tb03808.x}.

\noindent DAMODARAN, A. Equity risk premiums: determinants, estimation
and implications \textemdash{} the 2025 edition. \textbf{NYU Stern
School of Business working paper}, 2025. URL:
\url{https://pages.stern.nyu.edu/~adamodar/}.

\noindent DEMIGUEL, V.; GARLAPPI, L.; UPPAL, R. Optimal versus naive
diversification: how inefficient is the 1/N portfolio strategy?
\textbf{The Review of Financial Studies}, v. 22, n. 5, p. 1915-1953,
2009. DOI:
\href{https://doi.org/10.1093/rfs/hhm075}{10.1093/rfs/hhm075}.

\noindent FAMA, E. F. The behavior of stock-market prices.
\textbf{The Journal of Business}, v. 38, n. 1, p. 34-105, 1965.

\noindent FAMA, E. F.; FRENCH, K. R. Common risk factors in the returns
on stocks and bonds. \textbf{Journal of Financial Economics}, v. 33,
n. 1, p. 3-56, 1993.

\noindent HE, G.; LITTERMAN, R. The intuition behind Black-Litterman
model portfolios. \textbf{Goldman Sachs Investment Management Research
working paper}, 1999. URL:
\url{https://papers.ssrn.com/sol3/papers.cfm?abstract_id=334304}.

\noindent IDZOREK, T. M. A step-by-step guide to the Black-Litterman
model: incorporating user-specified confidence levels.
\textbf{Morningstar Inc working paper}, 2005. URL:
\url{https://corporate.morningstar.com/ib/documents/MethodologyDocuments/IBBAssociates/BlackLitterman.pdf}.

\noindent JOBSON, J. D.; KORKIE, B. M. Performance hypothesis testing
with the Sharpe and Treynor measures. \textbf{The Journal of Finance},
v. 36, n. 4, p. 889-908, 1981.

\noindent JORION, P. Bayes-Stein estimation for portfolio analysis.
\textbf{Journal of Financial and Quantitative Analysis}, v. 21, n. 3,
p. 279-292, 1986. DOI:
\href{https://doi.org/10.2307/2331042}{10.2307/2331042}.

\noindent KAHNEMAN, D. \textbf{Thinking, fast and slow}. New York:
Farrar, Straus and Giroux, 2011.
\href{https://us.macmillan.com/books/9780374533557/thinkingfastandslow}{Editor}.

\noindent KAHNEMAN, D. \textbf{Prêmio do Nobel de Ciências Econômicas
de 2002}. URL:
\url{https://www.nobelprize.org/prizes/economic-sciences/2002/kahneman/facts/}.

\noindent KAHNEMAN, D.; TVERSKY, A. Prospect theory: an analysis of
decision under risk. \textbf{Econometrica}, v. 47, n. 2, p. 263-291,
1979. DOI:
\href{https://doi.org/10.2307/1914185}{10.2307/1914185}.

\noindent KAN, R.; SMITH, D. R. The distribution of the sample
minimum-variance frontier. \textbf{Management Science}, v. 54, n. 7,
p. 1364-1380, 2008.

\noindent LEDOIT, O.; WOLF, M. Improved estimation of the covariance
matrix of stock returns with an application to portfolio selection.
\textbf{Journal of Empirical Finance}, v. 10, n. 5, p. 603-621, 2003.

\noindent LEDOIT, O.; WOLF, M. Honey, I shrunk the sample covariance
matrix. \textbf{The Journal of Portfolio Management}, v. 30, n. 4, p.
110-119, 2004. DOI:
\href{https://doi.org/10.3905/jpm.2004.110}{10.3905/jpm.2004.110}.

\noindent LÓPEZ DE PRADO, M. Building diversified portfolios that
outperform out of sample. \textbf{The Journal of Portfolio
Management}, v. 42, n. 4, p. 59-69, 2016.

\noindent MAILLARD, S.; RONCALLI, T.; TEÏLETCHE, J. The properties of
equally weighted risk contribution portfolios. \textbf{The Journal of
Portfolio Management}, v. 36, n. 4, p. 60-70, 2010. DOI:
\href{https://doi.org/10.3905/jpm.2010.36.4.060}{10.3905/jpm.2010.36.4.060}.

\noindent MALKIEL, B. G. \textbf{A random walk down Wall Street}.
W. W. Norton, 1973.
\href{https://wwnorton.com/books/9780393358384}{Editor}.

\noindent MARKOWITZ, H. Portfolio selection. \textbf{The Journal of
Finance}, v. 7, n. 1, p. 77-91, 1952. DOI:
\href{https://doi.org/10.2307/2975974}{10.2307/2975974}.

\noindent MARKOWITZ, H. \textbf{Prêmio do Nobel de Ciências Econômicas
de 1990}. URL:
\url{https://www.nobelprize.org/prizes/economic-sciences/1990/markowitz/facts/}.

\noindent MERTON, R. C. Lifetime portfolio selection under uncertainty:
the continuous-time case. \textbf{The Review of Economics and
Statistics}, v. 51, n. 3, p. 247-257, 1969.

\noindent MERTON, R. C. An analytic derivation of the efficient
portfolio frontier. \textbf{Journal of Financial and Quantitative
Analysis}, v. 7, n. 4, p. 1851-1872, 1972.

\noindent MERTON, R. C. An intertemporal capital asset pricing model.
\textbf{Econometrica}, v. 41, n. 5, p. 867-887, 1973.

\noindent MICHAUD, R. O. The Markowitz optimization enigma: is
``optimized'' optimal? \textbf{Financial Analysts Journal}, v. 45,
n. 1, p. 31-42, 1989. URL:
\url{https://www.jstor.org/stable/4479185}.

\noindent MICHAUD, R. O. \textbf{Efficient asset management: a
practical guide to stock portfolio optimization and asset allocation}.
Boston: Harvard Business School Press, 1998.

\noindent RONCALLI, T. \textbf{Introduction to risk parity and
budgeting}. Chapman \& Hall/CRC, 2013.
\href{https://www.crcpress.com/Introduction-to-Risk-Parity-and-Budgeting/Roncalli/p/book/9781482207156}{Editor}.

\noindent SHARPE, W. F. Capital asset prices: a theory of market
equilibrium under conditions of risk. \textbf{The Journal of Finance},
v. 19, n. 3, p. 425-442, 1964.

\noindent SHARPE, W. F. \textbf{Prêmio do Nobel de Ciências Econômicas
de 1990}. URL:
\url{https://www.nobelprize.org/prizes/economic-sciences/1990/sharpe/facts/}.

\noindent TOBIN, J. Liquidity preference as behavior towards risk.
\textbf{The Review of Economic Studies}, v. 25, n. 2, p. 65-86, 1958.

\noindent FAMA, E. F. \textbf{Prêmio do Nobel de Ciências Econômicas
de 2013}. URL:
\url{https://www.nobelprize.org/prizes/economic-sciences/2013/fama/facts/}.

\noindent CHALHOUB, L. \textbf{applied\_finance}: B3 market analytics
platform (Markowitz, $1/N$, Kahneman, Black-Litterman, ERC).
Repositório GitHub, 2026. URL:
\url{https://github.com/leonardochalhoub/applied_finance}.

\end{singlespace}

\end{document}
